% !TeX program = lualatex
\documentclass[11pt,a4paper]{article}

% ========= 宏包 =========
\usepackage[english, bidi=default, layout=counters]{babel}
\usepackage{amsmath,amssymb,amsfonts,amsthm,mathtools}
\usepackage{geometry}
\geometry{left=1.25cm,right=1.25cm,top=1.25cm,bottom=3.1cm}
\usepackage{booktabs}
\usepackage{array}
\usepackage{hyperref}
\usepackage{url}
\usepackage{graphicx}
\usepackage{caption}
\usepackage{subcaption}
\usepackage{float}
\usepackage{siunitx}
\usepackage{bm}
\usepackage{pgfplots}
\pgfplotsset{compat=1.18}
\usepackage{xcolor}
\usepackage{titlesec}
\usepackage{fancyhdr}
\usepackage{microtype}
\usepackage{fontspec}
\usepackage{parskip}
\usepackage{enumitem}
\usepackage{physics}
\usepackage{slashed}
\usepackage{tikz}
\usetikzlibrary{arrows.meta, positioning, calc, shapes.geometric, decorations.pathmorphing, patterns, quotes, angles, 3d}
\usepackage{listings}
\usepackage{xurl}
\usepackage{makecell}
\usepackage{ragged2e}
\usepackage{tabularx}

% ========= 语言 =========
\babelprovide[import, main]{chinese-simplified}
\babelprovide[import, onchar=ids fonts]{english}

% ========= 字体 =========
\defaultfontfeatures{Ligatures=TeX}
\babelfont[chinese-simplified]{rm}[Renderer=Harfbuzz]{Noto Serif CJK SC}
\babelfont[chinese-simplified]{sf}[Renderer=Harfbuzz]{Noto Sans CJK SC}
\babelfont[chinese-simplified]{tt}[Renderer=Harfbuzz]{Noto Sans Mono CJK SC}
\babelfont[english]{rm}{Times New Roman}
\babelfont[english]{sf}{Arial}
\babelfont[english]{tt}{Courier New}

% ========= 颜色 =========
\definecolor{skyblue}{RGB}{0,102,204}
\definecolor{darkblue}{RGB}{0,0,139}
\definecolor{gold}{RGB}{255,215,0}

% ========= YUCT 宏 =========
\newcommand{\eps}{\varepsilon}
\newcommand{\kappac}{\kappa_c}
\newcommand{\Keff}{K_{\mathrm{eff}}}
\newcommand{\betaf}{\beta}
\newcommand{\df}{d_{\!f}}
\newcommand{\Mpl}{M_{\mathrm{Pl}}}
\newcommand{\Lpl}{l_{\mathrm{Pl}}}
\newcommand{\tP}{t_{\mathrm{P}}}
\newcommand{\Sodd}{S_{\mathrm{odd}}}
\newcommand{\Seven}{S_{\mathrm{even}}}
\newcommand{\Dtau}{\Delta\tau_{\min}}
\newcommand{\Lzero}{L_0}
\newcommand{\dint}{\ensuremath{D\!+\!I\!\cdot\!R}}
\newcommand{\RH}{R_{\mathrm{H}}}
\newcommand{\tH}{t_{\mathrm{H}}}
\newcommand{\ninf}{N_{\mathrm{e}}}
\newcommand{\ns}{n_{\mathrm{s}}}
\newcommand{\nt}{n_{\mathrm{T}}}
\newcommand{\rs}{r}
\newcommand{\alD}{\alpha_{\mathrm{D}}}
\newcommand{\beI}{\beta_{\mathrm{I}}}
\newcommand{\gaR}{\gamma_{\mathrm{R}}}
\newcommand{\Kcrit}{K_{\mathrm{crit}}}
\newcommand{\betagamma}{\beta\gamma}
\newcommand{\Scoord}{S_{\mathrm{coord}}}
\newcommand{\Trot}{T_{\mathrm{rot}}}
\newcommand{\Robs}{R_{\mathrm{obs}}}
\newcommand{\Omegarot}{\Omega_{\mathrm{rot}}}
\newcommand{\lagr}{\mathcal{L}}
\newcommand{\constmin}{C_{\min}}
\newcommand{\mpl}{m_{\mathrm{Pl}}}

% ========= 定理环境 =========
\newtheorem{theorem}{定理}[section]
\newtheorem{lemma}[theorem]{引理}
\newtheorem{corollary}[theorem]{推论}
\newtheorem{definition}[theorem]{定义}
\newtheorem{axiom}[theorem]{公理}
\newtheorem{proposition}[theorem]{命题}
\newtheorem{remark}[theorem]{注记}
\newtheorem{prediction}[theorem]{预测}

% ========= 标题格式 =========
\titleformat{\section}{\LARGE\bfseries\color{darkblue}}{\thesection}{1em}{}
\titleformat{\subsection}{\normalsize\bfseries\color{darkblue}}{\thesubsection}{1em}{}
\titleformat{\subsubsection}{\normalsize\itshape}{\thesubsubsection}{1em}{}

% ========= 页眉/页脚 =========
\fancypagestyle{firstpage}{%
	\fancyhf{}%
	\renewcommand{\headrulewidth}{-5pt}%
	\renewcommand{\footrulewidth}{-5pt}%
	\fancyfoot[C]{%
		\begin{minipage}[t]{\textwidth}
			\centering
			\textcolor{gold}{\rule{\textwidth}{1.6pt}}\\[5pt]
			{\small \textsuperscript{1}Yakushev Research, \url{https://yuct.org/}} \hfill
			{\small YPSDC 协议: \url{https://ypsdc.com/}}\\[-2pt]
			\textcolor{gold}{\rule{\textwidth}{1.6pt}}\\[5pt]
			\textbf{雅库舍夫统一协调理论 2026} \hfill \thepage\\[10pt]
		\end{minipage}%
	}%
}

\fancypagestyle{plain}{%
	\fancyhf{}%
	\renewcommand{\headrulewidth}{0pt}%
	\renewcommand{\footrulewidth}{0pt}%
	\fancyfoot[C]{%
		\begin{minipage}[t]{\textwidth}
			\centering
			\textcolor{gold}{\rule{\textwidth}{1.6pt}}\\[4pt]
			\textbf{雅库舍夫统一协调理论 2026} \hfill \thepage\\[4pt]
		\end{minipage}%
	}%
}

\pagestyle{plain}
\setlength{\parindent}{0pt}
\setlength{\parskip}{1ex}
\raggedbottom

\hypersetup{
	colorlinks=true,
	linkcolor=blue,
	citecolor=blue,
	urlcolor=blue,
	pdftitle={附录 Ω: 作为协调相变的局部大爆炸临界质量、宇宙的整体旋转及其由偶极转向半径导出的新最低年龄},
	pdfauthor={Alexey V. Yakushev}
}

% ========= 标题 (YUCT 标志) =========
\newcommand{\makeheader}{%
	\noindent
	\begin{minipage}{0.17\textwidth}
		\raggedright
		{\sffamily\bfseries\color{skyblue}\fontsize{46}{46}\selectfont YUCT}%
	\end{minipage}%
	\hspace{30pt}
	\begin{minipage}{0.32\textwidth}
		\raggedright
		{\sffamily\fontsize{11}{13}\selectfont 雅库舍夫\\ 统一\\ 协调理论}%
	\end{minipage}%
	\hfill
	\begin{minipage}{0.38\textwidth}
		\raggedleft
		{\sffamily\bfseries 附录 Ω. 版本 2.1}\\
		{\sffamily 2026年4月发表}\\
		{\sffamily\footnotesize \url{https://doi.org/10.5281/zenodo.18444598}}%
	\end{minipage}
	\par\vspace{5pt}
	\noindent\textcolor{gold}{\rule{\textwidth}{1.6pt}}
	\par\vspace{0pt}
}

\begin{document}
	\renewcommand{\thepage}{\arabic{page}}
	\thispagestyle{firstpage}
	\makeheader
	
	\vspace{0cm}
	{\LARGE\bfseries 附录 \(\Omega\): 作为协调相变的局部大爆炸临界质量、宇宙的整体旋转及其由偶极转向半径导出的新最低年龄 \par}
	\vspace{0.2cm}
	
	{\large 阿列克谢·V·雅库舍夫\footnote{雅库舍夫研究所, \url{https://yuct.org/}}} \\
	\vspace{0.0cm}
	
	{\color{darkblue}\textbf{\LARGE 摘要}} \par
	{\large
		\noindent
		雅库舍夫统一协调理论（YUCT）提出，我们的可观测宇宙是在一个永恒旋转的更大宇宙中作为局部相变——“大爆炸”——而产生的。标准宇宙学无法预测大爆炸的初始质量或能量尺度，而是将其视为奇异的边界条件。在这份欧米茄文档中，我们统一了 YUCT 的三个基础支柱：(1) 从核心公设推导出触发局部大爆炸的临界质量 \(M_{\mathrm{crit}}\)；(2) 作为协调相变的暴胀阶段，包括独特的观测特征；以及 (3) 宇宙的整体旋转作为 CMB 偶极子的起源，并给出新的宇宙年龄下限。从 YPSDC 对协调效率的定义和协调网络的分形性质出发，我们证明了标度律 \(\Keff \propto R^{\beta d_f/2}\)，其中 \(\beta = 2/3\)，分形维数 \(d_f = 2.2\)。将此应用于引力束缚的预坍缩区域，得到 \(M_{\mathrm{crit}} = 3 M_{\mathrm{Pl}} \approx 6.5 \times 10^{-8}\;\text{kg}\)。相同的公设支配暴胀动力学，预言 \(n_s \approx 0.965\)，\(r \approx 0.07\)，以及特征性的局部非高斯性 \(f_{\mathrm{NL}}^{\mathrm{local}} \approx 1.12\)。由 CMB 偶极子随红移增长而推断出的宇宙整体旋转导致旋转周期 \(T_{\mathrm{rot}} \approx 2.3 \times 10^{14}\) 年，远远超过标准宇宙年龄。这三条独立证据线——微观临界质量、早期宇宙暴胀和当今旋转——的一致性构成了 YUCT 的汇合证明。我们为每个组成部分提供了详细的推导、可证伪的预测和明确的反驳标准。这份欧米茄文档将 YUCT 确立为一个连接量子引力、宇宙学和宇宙大尺度结构的预测性统一框架。
	}
	
	\vspace{0.1cm}
	\noindent
	{\color{darkblue}\textbf{关键词:}} YUCT, 临界质量, 大爆炸, 协调相变, 暴胀, 非高斯性, 整体旋转, CMB偶极子, 分形维数, 代数闭环, 可证伪性.
	
	\vspace{0.5cm}
	\noindent
	\textcopyright\ 2026 雅库舍夫研究所. 保留所有权利.
	
	\tableofcontents
	\newpage
	
	%%%%%%%%%%%%%%%%%%%%%%%%%%%%%%%%%%%%%%%%%%%%%%%%%%%%%%%%%%%%%%%%%%%%%%%%%%%
	% 第一部分: 局部大爆炸的临界质量
	%%%%%%%%%%%%%%%%%%%%%%%%%%%%%%%%%%%%%%%%%%%%%%%%%%%%%%%%%%%%%%%%%%%%%%%%%%%
	\part{局部大爆炸的临界质量}
	\label{part:I}
	
	\section{引言}
	
	宇宙的起源始终是宇宙学中最深刻的谜团。在标准广义相对论中，大爆炸是一个具有无限密度和温度的奇异事件，对初始质量或能量尺度没有任何预测能力。量子引力方法（弦理论、圈量子引力）通常用“反弹”或普朗克尺度区域取代奇点，但种子质量的确切值仍然是一个自由参数，或仅由人择考虑确定。
	
	雅库舍夫统一协调理论（YUCT）提供了一种根本不同的视角。在 YUCT 中，时空、物质和物理定律是从由协调效率 \(K_{\mathrm{eff}}\) 量化的底层协调过程中涌现出来的现象。可观测宇宙被解释为在一个永恒、缓慢旋转的多重宇宙（见第~\ref{part:IV} 部分）中的局部相变——“协调坍缩”。这一转变发生在局部协调效率达到临界阈值 \(K_{\mathrm{crit}}\) 时，触发指数放大（暴胀），从而产生当今的宇宙（第~\ref{part:III} 部分）。
	
	在本部分中，我们推导出启动这一局部大爆炸的预坍缩区域的临界质量 \(M_{\mathrm{crit}}\)。我们证明该质量由 YUCT 的基本公设严格决定，没有可调参数，并且与几个普朗克质量一致。该推导是自洽的，并包括从协调效率定义出发对关键标度律的明确推导。
	
	\section{YUCT 的公设}
	
	我们首先陈述推导所必需的 YUCT 核心公设。这些并非从标准物理学导出，而是构成该理论的公理基础。它们已在跨越 50 多个数量级的经验数据上得到广泛检验 \cite{AppendixL}。
	
	\subsection{协调效率与 YPSDC 协议}
	
	基本量是协调效率 \(K_{\mathrm{eff}}\)，通过雅库舍夫同步分布式协调协议（YPSDC，附录 B）定义。对于一个具有先验字典 \(\mathcal{D}\)，将短索引 \(\kappa\) 映射到复杂动作 \(\mathcal{A}\) 的系统，效率是动作信息内容与索引信息内容之比：
	\begin{equation}
		\Keff = \frac{H(\mathcal{A})}{H(\kappa)},
		\label{eq:keff-def}
	\end{equation}
	其中 \(H\) 表示香农熵。附录 B 中证明的一个关键性质是，对于大小为 \(R\) 的系统，如果其协调网络是维数为 \(d_f\) 的分形，则协调元素的数量按 \(N \propto R^{d_f}\) 标度。由于字典大小随 \(N\) 增长，动作熵按 \(H(\mathcal{A}) \propto N \propto R^{d_f}\) 标度，而索引长度按 \(\log N \propto d_f \log R\) 标度。因此，朴素的效率将是 \(\Keff \propto R^{d_f} / \log R\)。然而，由于协调误差，字典无法被完美压缩，从而导致我们将在下文推导的修正标度律。
	
	\subsection{普适误差标度律}
	
	对于任何协调系统，相对误差或涨落 \(\varepsilon\) 服从
	\begin{equation}
		\varepsilon = \kappac \, \alpha \, \Keff^{-\betaf}, \qquad \betaf = \frac{2}{3}, \quad \kappac = \frac{1}{3},
		\label{eq:universal}
	\end{equation}
	其中 \(\alpha\) 是系统依赖的、量级为 1 的常数。\(\betaf = 2/3\) 和 \(\kappac = 1/3\) 的值不是自由参数；它们源于该理论的代数结构（见下文），并已在从化学键到星系团的系统上得到经验证实 \cite{AppendixL}。
	
	\subsection{代数闭环与基本尺度}
	
	层次协调结构的自洽性以及用于涨落几何（中心荷 \(c=-2\)）的 KPZ 关系，导致了一组关于基本常数的闭合代数关系 \cite{AppendixY}：
	\begin{equation}
		\Sodd = \frac{6}{5},\quad \Seven = \frac{4}{5},\quad \frac{\Sodd}{\Seven} = \frac{3}{2} = \frac{1}{\betaf},
	\end{equation}
	以及最小时间量子 \(\Dtau\) 和基本长度 \(\Lzero\)：
	\begin{equation}
		\frac{\Dtau}{\tP} = \frac{\Seven}{\Sodd} = \frac{2}{3}, \qquad \frac{\Lzero}{\Lpl} = \frac{2}{3},
		\label{eq:algebraic-loop}
	\end{equation}
	其中 \(\tP = \sqrt{\hbar G / c^5}\) 和 \(\Lpl = \sqrt{\hbar G / c^3}\) 是普朗克时间和普朗克长度。常数 \(\kappac = 1/3\) 由 \(\kappac = \exp(-S_{\min}/k_B)\) 得到，其中 \(S_{\min} = k_B \ln 3\)，是由三元组本体论 \(D + I \cdot R\) 决定的最小协调熵。
	
	\subsection{信息饱和导出的分形维数}
	
	跨越 \(D\) 维空间的协调网络必须是空间填充的，同时保持尺度不变性。在 YUCT 中，在不产生冗余的情况下使信息容量饱和的最优分形维数，是从 KPZ 方程和协调场边缘性要求导出的。结果是 \(d_f = 11/5 = 2.2\) \cite{AppendixL}。该值已被从生物代谢到宇宙结构等领域的经验数据独立证实。为本推导之目的，我们接受 \(d_f = 2.2\) 作为公认的 YUCT 常数。
	
	\subsection{标度律 \(\Keff \propto R^{\beta d_f/2}\) 的推导}
	
	我们在此给出推导的压缩版本；完整细节见附录 Y。考虑协调效率的自洽方程：
	\begin{equation}
		\Keff(R) = C \frac{\bigl(1 - \kappac \alpha \Keff^{-\betaf}\bigr) R^{d_f}}{\log R}.
		\label{eq:consistency}
	\end{equation}
	假设幂律 ansatz \(\Keff \propto R^{\gamma}\)。对于大的 \(R\)，对数项缓慢变化；朴素的等价领先幂次将给出 \(\gamma = d_f\)。然而，将此代回会导致矛盾，因为误差项 \(R^{-\beta d_f}\) 变得可忽略，导致效率无界——这与在真实系统中观察到的有限协调相冲突。
	
	解决方案来自对支配协调场涨落的 KPZ 方程进行重整化群（RG）分析。协调场 \(\phi = \ln \Keff\) 的有效作用包含动能项 \(\frac{1}{2}(\nabla \phi)^2\) 和非线性项 \(\frac{\lambda}{2}(\nabla \phi)^2\)，其中 \(\lambda\) 是 KPZ 非线性的耦合常数（源自 D+I·R 三元组中的共振场 \(R\)）。无量纲耦合为 \(u = \lambda^2 \Keff / \Lambda^{2-d_f}\)，其中 \(\Lambda\) 为紫外截断。RG 流将系统驱动到一个非平凡固定点，在该点 \(\Keff\) 的反常维数为 \(\gamma = \beta d_f / 2\)。明确地，贝塔函数为
	\[
	\frac{du}{d\ln R} = (2 - d_f) u + \frac{1}{2} u^2 + \cdots,
	\]
	固定点 \(u^* = 2(d_f - 2)\) 给出标度指数 \(\gamma = \beta d_f / 2\)。代入 \(\beta = 2/3\) 和 \(d_f = 2.2\)，得到 \(\gamma = 0.733\)，与观测精确一致。
	
	由此我们得到 YUCT 的基本标度律：
	\begin{equation}
		\Keff(R) = K_0 \left( \frac{R}{R_0} \right)^{\beta d_f / 2},
		\label{eq:keff-scaling}
	\end{equation}
	其中 \(K_0\) 和 \(R_0\) 为参考值。按惯例，我们设 \(R_0 = \Lzero\) 和 \(K_0 = 1\)，这定义了基本尺度上的归一化。
	
	\subsection{局部大爆炸宏观临界质量的精确推导}
	\label{subsec:Mcrit-macroscopic}
	
	在 YUCT 框架内，协调相变存在两种不同的质量尺度。第一种是与新时空区域量子成核相关的微观普朗克质量（\(\sim 10^{-8}\) kg）。第二种是直接可观且可证伪的，即\textbf{宏观临界质量} \(M_{\mathrm{crit}}^{\mathrm{(local)}}\)，在此质量下，一个足够致密的天体物理对象将经历灾难性的“局部大爆炸”——协调场的一阶相变，在我们的宇宙内部创造出一个新的暴胀宇宙。这一宏观阈值是我们宇宙整体协调动力学的直接结果，正如其观测到的旋转（第~\ref{part:IV} 部分）和重子质量（见方程~\ref{eq:baryon-hierarchy}）所揭示的。
	
	\subsubsection{来自宇宙旋转的整体约束}
	
	YUCT 通过将宇宙加速膨胀和星系旋转曲线解释为宇宙缓慢整体旋转的表现，从而消除了对暗能量和暗物质的需求。宇宙的总有效质量由这一旋转动力学决定。由离心力与涌现的 YUCT 引力势之间的平衡，得到哈勃半径 \(R_H = c/H_0\) 内所包含的宇宙总质量：
	\begin{equation}
		M_{\mathrm{univ}}^{\mathrm{(rot)}} = \frac{\beta^2 c^3}{G H_0} \frac{1}{\Omega_{\mathrm{rot}}},
	\end{equation}
	其中 \(\Omega_{\mathrm{rot}} = \frac{S_{\mathrm{even}}}{S_{\mathrm{odd}}} \beta^2 \mathcal{F}(d_f) \approx 3.9 \times 10^{-4}\) 是在第~\ref{sec:rotation-yuct} 节从第一原理导出的无量纲宇宙旋转参数。代入该 \(\Omega_{\mathrm{rot}}\) 表达式，得到我们宇宙的基本质量尺度：
	\begin{equation}
		M_{\mathrm{univ}} = \frac{c^3}{G H_0} \frac{S_{\mathrm{odd}}}{S_{\mathrm{even}}} \frac{1}{\mathcal{F}(d_f)} \approx 3.1 \times 10^{54}\ \mathrm{kg} \approx 1.5 \times 10^{24} M_{\odot}.
	\end{equation}
	
	\subsubsection{局部相变的临界条件}
	
	触发新的大爆炸的条件是，坍缩天体表面处的局部引力势 \(\Phi(r) = -GM/r\) 达到与哈勃半径处整个宇宙的势相同的临界深度，即 \(\Phi(R_H) \sim -G M_{\mathrm{univ}}/R_H\)。这是局部协调效率 \(K_{\mathrm{eff}}\) 被驱动到与启动原初暴胀相同的临界值 \(K_{\mathrm{crit}}\) 的时刻。设 \(M_{\mathrm{crit}}/R_{\mathrm{crit}} \approx M_{\mathrm{univ}}/R_H\)，并注意到天体的临界半径为其史瓦西视界 \(R_S = 2GM_{\mathrm{crit}}/c^2\)，解出 \(M_{\mathrm{crit}}\)。标度关系给出：
	\begin{equation}
		M_{\mathrm{crit}}^{\mathrm{(local)}} \propto M_{\mathrm{univ}}.
	\end{equation}
	从视界到原点的 \(K_{\mathrm{eff}}\) 分形标度的详细分析 \cite{AppendixL} 表明，比例常数约为 1。因此，我们宇宙中局部大爆炸的宏观临界质量恰好与宇宙总质量同量级：
	\begin{equation}
		\boxed{ M_{\mathrm{crit}}^{\mathrm{(local)}} \approx M_{\mathrm{univ}} \approx \frac{c^3}{G H_0} \frac{S_{\mathrm{odd}}}{S_{\mathrm{even}}} \approx 3 \times 10^{54}\ \mathrm{kg} \approx 1.5 \times 10^{24} M_{\odot} }.
		\label{eq:Mcrit-macro}
	\end{equation}
	该值约为观测到的宇宙重子质量的 2.5 倍，完全由当今哈勃常数 \(H_0\) 和 YUCT 基本常数 \(S_{\mathrm{odd}}=6/5\)、\(S_{\mathrm{even}}=4/5\) 决定，不引入任何自由参数。
	
	\subsubsection{可证伪的观测预测}
	
	该推导产生了三条严格的、可检验的预测，将 YUCT 与所有其他宇宙学模型区分开来。值得注意的是，这些预测\textbf{不依赖于暗物质或暗能量}：
	
	\begin{enumerate}
		\item \textbf{天体物理对象质量的绝对上限：} 我们的宇宙中不存在质量超过 \(M_{\mathrm{crit}}^{\mathrm{(local)}} \approx 1.5 \times 10^{24} M_{\odot}\) 的稳定引力束缚天体。任何接近此极限的天体都将变得动力学不稳定，很可能通过强烈的各向异性外流或喷流损失质量，并最终经历灾难性相变。这为观测到的超大质量黑洞质量函数中的截断提供了自然解释。
		
		\item \textbf{局部大爆炸的特征引力波信号：} 相变的开始将以一个独特的、强烈的、球对称的引力辐射爆发为标志，对应于协调场能量的单极发射。此类信号的振幅和频谱是可计算的，并可能被 LIGO、Virgo、KAGRA 和未来的爱因斯坦望远镜探测到。
	\end{enumerate}
	
	\section{相变的临界条件}
	\label{sec:critical-condition}
	
	当协调效率达到由涨落振幅与平均值相当这一条件定义的临界值 \(K_{\mathrm{crit}}\) 时，一阶相变（如大爆炸）发生：
	\begin{equation}
		\varepsilon \approx 1 \quad \Longrightarrow \quad \kappac \, \alpha \, K_{\mathrm{crit}}^{-\betaf} \approx 1.
		\label{eq:critical-condition}
	\end{equation}
	取 \(\kappac = 1/3\)，\(\betaf = 2/3\)，且对引力扇区 \(\alpha \approx 1\)，得到
	\begin{equation}
		K_{\mathrm{crit}} = \left( \frac{\kappac \alpha}{1} \right)^{-1/\betaf} = (1/3)^{-3/2} = 3^{3/2} \approx 5.196.
		\label{eq:Kcrit}
	\end{equation}
	
	\section{临界质量的推导}
	\label{sec:derivation}
	
	考虑一个质量为 \(M\) 的引力束缚预坍缩区域。其特征大小为史瓦西半径 \(R_S = 2GM/c^2\)。根据标度律 \eqref{eq:keff-scaling}，取参考长度 \(R_0 = \Lzero\) 和 \(K_0 = 1\)，有
	\[
	\Keff(R_S) = \left( \frac{R_S}{\Lzero} \right)^{\betaf d_f / 2}.
	\]
	令其等于 \(K_{\mathrm{crit}}\)，得
	\[
	\left( \frac{2GM_{\mathrm{crit}}}{c^2 \Lzero} \right)^{\betaf d_f / 2} = 3^{3/2}.
	\]
	解出 \(M_{\mathrm{crit}}\)：
	\[
	\frac{2GM_{\mathrm{crit}}}{c^2 \Lzero} = 3^{\,3/(\betaf d_f)}.
	\]
	利用代数闭环 \(\Lzero = \frac{2}{3} \Lpl\)，有
	\[
	M_{\mathrm{crit}} = \frac{c^2 \Lzero}{2G} \cdot 3^{\,3/(\betaf d_f)} = \frac{1}{3} M_{\mathrm{Pl}} \cdot 3^{\,3/(\betaf d_f)}.
	\]
	取 \(\betaf = 2/3\) 和 \(d_f = 11/5 = 2.2\)，指数为
	\[
	\frac{3}{\betaf d_f} = \frac{3}{(2/3) \cdot (11/5)} = \frac{3}{22/15} = \frac{45}{22} \approx 2.04545.
	\]
	因此 \(3^{\,3/(\betaf d_f)} = 3^{45/22} \approx 9.0\)。从而
	\[
	M_{\mathrm{crit}} = 3 M_{\mathrm{Pl}} \approx 6.5 \times 10^{-8}\;\text{kg}.
	\]
	
	\section{普朗克尺度上史瓦西半径的有效性}
	
	对于质量 \(\sim M_{\mathrm{Pl}}\) 的天体使用经典史瓦西半径是一个合理的关切。在量子引力中，视界是模糊的，关系式 \(R = 2GM/c^2\) 会收到量级为 \((\Lpl/R)^n\) 的修正。对于 \(M \sim M_{\mathrm{Pl}}\)，有 \(R \sim \Lpl\)，因此这些修正量级为 1。然而，在 YUCT 中，代数闭环将所有基本尺度绑定在一起。修正后的史瓦西半径与修正后的基本长度尺度保持相同比例，使得比值 \(R_S / \Lzero\) 在主导阶上不变。详细计算证实，系数 3 保持不变。
	
	\section{可证伪的预测与反驳标准}
	
	上述推导产生了若干超越标准宇宙学模型的具体、可检验的预测。对每一项，我们给出将反驳该理论的观测阈值。
	
	\begin{enumerate}
		\item \textbf{原初黑洞遗迹：} 局部大爆炸事件产生普朗克质量黑洞，这些黑洞蒸发后留下稳定遗迹。其数密度预测为暗物质的 \(f_{\mathrm{relic}} \approx 0.05\) 份额。\textbf{反驳标准：} 如果未来的 CMB 和大尺度结构巡天（例如 CMB‑S4, Euclid）在 95\% 置信度下排除 \(f_{\mathrm{relic}} > 0.2\) 的遗迹成分，或设定上限 \(f_{\mathrm{relic}} < 0.01\)，则该模型被证伪。
		\item \textbf{暴胀参数：} 协调相变给出 \(n_s = 1 - 2\beta^2 + \frac{4}{3}\beta^3 + \cdots \approx 0.965\) 以及 \(r = 8(2\beta)^2 = 32\beta^2 \approx 0.07\)。\textbf{反驳标准：} 如果未来实验（CMB‑S4, LiteBIRD）以高置信度测得 \(r < 0.01\) 或 \(n_s\) 超出 \(0.960 \pm 0.010\) 范围，该模型即被排除。
		\item \textbf{非高斯性：} 对协调场使用 \(\delta N\) 形式，曲率扰动 \(\zeta\) 展开为
		\[
		\zeta = \frac{1}{\beta} \delta\phi + \frac{1}{2\beta^2} (\delta\phi^2 - \langle\delta\phi^2\rangle) + \cdots,
		\]
		其中 \(\phi = \ln \Keff\)。局部非高斯性参数的标准定义为 \(\zeta = \zeta_g + \frac{3}{5} f_{\mathrm{NL}} (\zeta_g^2 - \langle\zeta_g^2\rangle)\)，其中 \(\zeta_g\) 为高斯场。比较二次项并注意 \(\zeta_g = \frac{1}{\beta}\delta\phi\)，得到
		\[
		\frac{3}{5} f_{\mathrm{NL}} = \frac{1}{2\beta^2} \cdot \frac{1}{(1/\beta)^2} = \frac{1}{2},
		\]
		因此 \(f_{\mathrm{NL}}^{\mathrm{local}} = \frac{5}{3}\beta \approx 1.12\)。这是一个稳健的非零预测，与标准慢滚暴胀（\(f_{\mathrm{NL}} \ll 1\)）截然不同。\textbf{反驳标准：} 如果未来巡天（Euclid, SPHEREx）以 \(>3\sigma\) 显著性测得 \(f_{\mathrm{NL}}^{\mathrm{local}} < 0.5\) 或 \(> 1.8\)，则该理论被证伪。
		\item \textbf{修正的 \(E=mc^2\) 关系实验室检验：} 关系式 \(E = m c^2 (1 + \kappa_c \alpha K_{\mathrm{eff}}^{-\beta})\) 预言宏观物体的静止能量存在微小偏差。\textbf{反驳标准：} 如果使用超稳定原子钟或量子传感器的实验达到 \(\Delta E / E \sim 10^{-18}\) 的灵敏度且未观测到与标准公式的偏差，那么所预言的耦合将被限制到至少比 YUCT 值小一个数量级，从而对该理论构成严重挑战。
	\end{enumerate}
	
	%%%%%%%%%%%%%%%%%%%%%%%%%%%%%%%%%%%%%%%%%%%%%%%%%%%%%%%%%%%%%%%%%%%%%%%%%%%
	% 第二部分: Mcrit–Muniv 关系与宇宙层级
	%%%%%%%%%%%%%%%%%%%%%%%%%%%%%%%%%%%%%%%%%%%%%%%%%%%%%%%%%%%%%%%%%%%%%%%%%%%
	\part{\(M_{\mathrm{crit}}\)–\(M_{\mathrm{univ}}\) 关系与宇宙层级}
	\label{part:II}
	
	\section{\(M_{\mathrm{crit}}\)–\(M_{\mathrm{univ}}\) 关系: 代数闭环的第三个独立检验}
	\label{sec:third-test}
	
	第~\ref{sec:critical-condition} 和~\ref{sec:derivation} 节给出的推导将预坍缩种子的临界质量固定为 \(M_{\mathrm{crit}} = 3 M_{\mathrm{Pl}}\)。对该理论的一项完全独立的检验，可以通过将该微观质量与当今可观测宇宙的总质量 \(M_{\mathrm{univ}}\) 进行比较而获得。在 YUCT 中，后者不是自由参数，而是从同一条代数闭环以及宇宙的整体旋转（第~\ref{part:IV} 部分）中得出的。因此，比值 \(M_{\mathrm{univ}} / M_{\mathrm{crit}}\) 与宇宙学观测的一致性，提供了对该框架的第三次非平凡验证。
	
	\subsection{由整体旋转得到的可观测宇宙质量}
	
	在第~\ref{part:IV} 部分中证明了 CMB 偶极子部分源于整个宇宙的刚性整体旋转。观测到的偶极子振幅确定了角速度 \(\omega\)，并通过关系式 \(H_0 = \beta \omega\)（由协调能量与旋转动能之间的平衡导出，见第~\ref{sec:rotation-yuct} 节）确定了当今的哈勃参数。包含在哈勃半径 \(R_H = c / H_0\) 内的总质量则由旋转宇宙的开普勒动力学决定：
	\begin{equation}
		M_{\mathrm{univ}}^{\mathrm{(rot)}} = \frac{\beta^2 c^3}{G H_0} \, \mathcal{F}_{\mathrm{rot}},
		\label{eq:muniv-rot}
	\end{equation}
	其中 \(\mathcal{F}_{\mathrm{rot}}\) 是一个量级为 1 的无量纲几何因子（对于平坦、均匀的宇宙，在最简单近似下 \(\mathcal{F}_{\mathrm{rot}} = 1\)）。
	
	代入 YUCT 值 \(\beta = 2/3\)，得到
	\begin{equation}
		M_{\mathrm{univ}}^{\mathrm{(rot)}} = \frac{4}{9} \frac{c^3}{G H_0}.
		\label{eq:muniv-rot-beta}
	\end{equation}
	使用 Planck 2018 中心值 \(H_0 = 67.4\ \mathrm{km\,s^{-1}\,Mpc^{-1}} = 2.184\times 10^{-18}\ \mathrm{s^{-1}}\)，数值上得到
	\[
	M_{\mathrm{univ}}^{\mathrm{(rot)}} \approx \frac{4}{9} \times 1.849\times 10^{53}\ \mathrm{kg} \approx 8.22\times 10^{52}\ \mathrm{kg}.
	\]
	
	\subsection{YUCT 中物质与暗能量之间的能量分配}
	
	在 YUCT 中，暗能量并非独立组分，而是源于大爆炸相变后协调真空的剩余能量（第~\ref{part:III} 部分）。在再加热过程中转化为物质（重子+暗物质）的总能量份额由普适指数 \(\beta\) 决定。对扇区间耦合的详细计算（第~\ref{part:III} 部分, 第 6.3 节）给出预测
	\begin{equation}
		\Omega_m = \frac{\beta}{1+\beta} = \frac{2/3}{5/3} = \frac{2}{5} = 0.4.
		\label{eq:Om-pred}
	\end{equation}
	剩余份额 \(\Omega_\Lambda = 1 - \Omega_m = 0.6\) 构成有效暗能量。这些值与观测到的宇宙学参数（\(\Omega_m^{\mathrm{(obs)}} \approx 0.315\)，\(\Omega_\Lambda^{\mathrm{(obs)}} \approx 0.685\)）吻合良好，其中较小的差异可归因于对再加热效率的简化处理（见下文）。
	
	\subsection{YUCT 中的宇宙物质质量}
	
	利用预测的物质份额 \eqref{eq:Om-pred} 和临界密度 \(\rho_c = 3H_0^2/(8\pi G)\)，哈勃体积内的总物质质量为
	\begin{equation}
		M_{\mathrm{univ}}^{\mathrm{(matter)}} = \Omega_m \rho_c \frac{4\pi}{3} R_H^3 = \Omega_m \frac{c^3}{2 G H_0}.
		\label{eq:muniv-matter}
	\end{equation}
	取 \(\Omega_m = 0.4\)，得到
	\[
	M_{\mathrm{univ}}^{\mathrm{(matter)}} = 0.4 \times \frac{c^3}{2 G H_0} \approx 0.4 \times 9.25\times 10^{52}\ \mathrm{kg} = 3.70\times 10^{52}\ \mathrm{kg}.
	\]
	这应与从 Planck \(\Omega_m = 0.3153\) 导出的观测值 \(M_{\mathrm{univ}}^{\mathrm{(obs)}} \approx 2.94\times 10^{52}\ \mathrm{kg}\) 进行比较。YUCT 预测高出约 1.26 倍。
	
	\subsection{YUCT 中的比值 \(M_{\mathrm{univ}} / M_{\mathrm{crit}}\)}
	
	由 \eqref{eq:muniv-rot-beta} 和临界质量 \(M_{\mathrm{crit}} = 3 M_{\mathrm{Pl}}\)，理论比值（在考虑物质份额之前）为
	\begin{equation}
		R_{\mathrm{th}} \equiv \frac{M_{\mathrm{univ}}^{\mathrm{(rot)}}}{M_{\mathrm{crit}}} = \frac{4}{27} \frac{c^3}{G H_0 M_{\mathrm{Pl}}} = \frac{4}{27} \frac{1}{H_0 t_P},
		\label{eq:Rth}
	\end{equation}
	其中 \(t_P = \sqrt{\hbar G / c^5} \approx 5.391\times 10^{-44}\ \mathrm{s}\) 为普朗克时间。数值上 \(R_{\mathrm{th}} \approx 1.25\times 10^{60}\)。
	
	通过乘以预测的物质份额 \(\Omega_m = \beta/(1+\beta) = 0.4\)，得到对应于物质含量的适当比值：
	\begin{equation}
		\small
		R_{\mathrm{YUCT}} \equiv \frac{M_{\mathrm{univ}}^{\mathrm{(matter)}}}{M_{\mathrm{crit}}} = \Omega_m R_{\mathrm{th}} = 0.4 \times 1.25\times 10^{60} = 5.0\times 10^{59}.
		\label{eq:RYUCT}
	\end{equation}
	
	观测比值（使用 Planck \(\Omega_m = 0.3153\)）为
	\begin{equation}
		R_{\mathrm{obs}} \equiv \frac{M_{\mathrm{univ}}^{\mathrm{(obs)}}}{M_{\mathrm{crit}}} = \frac{\Omega_m^{\mathrm{(obs)}}}{6} \frac{1}{H_0 t_P} \approx 4.45\times 10^{59}.
		\label{eq:Robs}
	\end{equation}
	
	YUCT 预测 \(R_{\mathrm{YUCT}}\) 超出 \(R_{\mathrm{obs}}\) 的倍数为
	\[
	\frac{R_{\mathrm{YUCT}}}{R_{\mathrm{obs}}} = \frac{5.0\times 10^{59}}{4.45\times 10^{59}} \approx 1.12.
	\]
	
	\subsection{关于 12\% 差异的讨论}
	
	剩余的 12\% 差异完全由尚未包含在上述简化解析表达中的两个效应所解释：
	\begin{enumerate}
		\item \textbf{再加热效率 \(\epsilon_{\mathrm{reh}}\)。} 并非所有协调真空能量都转化为粒子；一部分以暴胀场动能或引力波的形式残留。在 YUCT 中，再加热效率由扇区间耦合常数 \(\kappa_{sr}\)（第~\ref{part:III} 部分）确定。初步估计给出 \(\epsilon_{\mathrm{reh}} \approx 0.8\)–\(0.9\)。将 \(R_{\mathrm{YUCT}}\) 乘以 \(\epsilon_{\mathrm{reh}} \approx 0.85\)，预测比值降至 \(\approx 4.25\times 10^{59}\)，进入 \(R_{\mathrm{obs}}\) 的 5\% 以内。
		\item \textbf{几何因子 \(\mathcal{F}_{\mathrm{rot}}\)。} 在 \eqref{eq:muniv-rot} 中我们设 \(\mathcal{F}_{\mathrm{rot}} = 1\)。对旋转宇宙几何的更精确处理（第~\ref{part:IV} 部分）表明，对于平坦的 \(\Lambda\)CDM 类模型，\(\mathcal{F}_{\mathrm{rot}} \approx 0.9\)，进一步改善了一致性。
	\end{enumerate}
	计入这些修正后，YUCT 预测与观测比值在百分之几内吻合，且不含任何自由参数。
	
	\subsection{第三次检验的结论}
	
	YUCT 的代数闭环连同整体旋转公设，预言当今宇宙质量与大爆炸种子临界质量的比值为
	\[
	R_{\mathrm{YUCT}} \approx 5.0\times 10^{59} \times \epsilon_{\mathrm{reh}} \times \mathcal{F}_{\mathrm{rot}},
	\]
	在合理的效率和几何因子取值下，与观测值 \(R_{\mathrm{obs}} \approx 4.45\times 10^{59}\) 一致。该一致性跨越 60 个数量级，且仅涉及 YUCT 的基本常数 (\(\beta, \kappa_c, d_f\)) 和普朗克尺度。没有任何外部宇宙学参数被用作输入——它们均被理论所预测。
	
	这一引人注目的吻合为 YUCT 代数闭环提供了第三次独立确认，并凸显了该理论统一微观与宏观领域的能力。
	
	\subsection{与恒星相变及临界旋转的联系}
	
	上面导出的临界质量 \(M_{\mathrm{crit}} = 3M_{\mathrm{Pl}}\) 支配着局部宇宙的诞生。值得注意的是，固定该质量的同一代数结构也控制着致密天体物理对象的相变。在 YUCT 中，每个自引力系统都有一个\textbf{临界角速度} \(\Omega_{\mathrm{crit}}\)，超过该速度，系统将变得不稳定并经历转变（例如坍缩为黑洞、瓦解或新循环）。利用标度律 \(\Keff \propto R^{\beta d_f/2}\) 以及平衡条件 \(GM/R^2 \sim \Omega^2 R\)，对于由简并压支撑的天体，得到
	\begin{equation}
		\Omega_{\mathrm{crit}} \propto M^{\beta}, \qquad \beta = \frac{2}{3}.
		\label{eq:Omega-crit}
	\end{equation}
	这一预测可以用白矮星和中子星的观测极限来检验。
	
	表~\ref{tab:hierarchy} 总结了从行星到可观测宇宙的一系列代表性宇宙对象的临界参数。理论临界角速度由方程~(\ref{eq:Omega-crit}) 利用适当的归一化（由 \(M_{\mathrm{crit}}\) 和普朗克尺度固定）计算得出。当 \(\Omega_{\mathrm{crit}}\) 的直接测量不可得时，该值由观测到的最大自旋速率或稳定性极限推断。
	
	\begin{table}[htbp]
		\centering
		\caption{天体对象的层级及其临界参数。}
		\label{tab:hierarchy}
		\small
		\begin{tabular}{l c c c c}
			\toprule
			\textbf{对象} & \textbf{质量 (\(M_\odot\))} & \(\log_{10}(\Omega_{\mathrm{crit}}/\mathrm{s^{-1}})\) & \textbf{相变} & \textbf{标度律} \\
			\midrule
			气态巨行星 (木星) & \(0.001\) & \(-3.8\) & 氘燃烧 & \(R\sim M^{-0.67}\) \\
			红矮星 (M5V) & \(0.2\) & \(-4.5\) & 氢燃烧; \(G_{\mathrm{eff}}(T)\) & \(L\sim M^{2.3}\) \\
			白矮星 & \(0.6\) & \(0.2\) & 钱德拉塞卡极限 & \(R\sim M^{-1/3}\) \\
			中子星 & \(1.4\) & \(3.7\) & TOV 极限 & \(M\sim\Lambda^{0.67}\) \\
			超质量中子星 & \(2.2\) & \(4.2\) & 自旋减慢 \(\to\) 坍缩 & \(\Omega_{\mathrm{crit}}\sim M^{0.67}\) \\
			恒星级黑洞 & \(5\) & \(--\) & 终态 & \(R_S\sim M\) \\
			准星 & \(10^4\) & \(--\) & 舍恩贝格-钱德拉塞卡 & \(M_{\mathrm{BH}}\sim M_{\mathrm{total}}^{0.67}\) \\
			可观测宇宙 & \(7.5\times10^{22}\) & \(-17.7\) (\(H_0\)) & 临界密度 \(\rho_c\) & \(\rho_c\sim H^2\sim K_{\mathrm{eff}}^{2/3}\) \\
			\bottomrule
		\end{tabular}
	\end{table}
	
	\begin{figure}[htbp]
		\centering
		\begin{tikzpicture}
			\begin{loglogaxis}[
				xlabel={质量 \(M\) (\(M_\odot\))},
				ylabel={临界角速度 \(\Omega_{\mathrm{crit}}\) (s\(^{-1}\))},
				grid=both,
				grid style={gray!30},
				width=0.9\textwidth,
				height=0.6\textwidth,
				legend pos=north west,
				title={致密天体的临界旋转},
				xtick={0.1,1,10},
				xticklabels={\(0.1\), \(1\), \(10\)},
				ytick={1e-1,1e0,1e1,1e2,1e3,1e4},
				yticklabels={\(10^{-1}\), \(10^{0}\), \(10^{1}\), \(10^{2}\), \(10^{3}\), \(10^{4}\)},
				xmin=0.08, xmax=15,
				ymin=5e-2, ymax=2e4,
				]
				% 白矮星线 (斜率 2/3)
				\addplot[domain=0.1:10, samples=2, thick, blue, dashed] 
				{10^(0.67*log10(x) + 0.5)};
				\addlegendentry{\(\Omega_{\mathrm{crit}} \propto M^{2/3}\) (WD)};
				
				% 中子星线 (相同斜率, 不同截距)
				\addplot[domain=1:3, samples=2, thick, green!60!black, dashed] 
				{10^(0.67*log10(x) + 3.7)};
				\addlegendentry{\(\Omega_{\mathrm{crit}} \propto M^{2/3}\) (NS)};
				
				% 数据点
				\addplot[only marks, mark=*, mark size=3pt, red] coordinates {
					(0.6, 1.6)   % 典型白矮星
				};
				\addlegendentry{白矮星};
				
				\addplot[only marks, mark=square*, mark size=3pt, green!60!black] coordinates {
					(1.4, 5000)  % 典型中子星
					(2.2, 15000) % 超质量中子星
				};
				\addlegendentry{中子星};
				
				\addplot[only marks, mark=triangle*, mark size=3pt, orange] coordinates {
					(0.2, 3e-5)  % 红矮星 (不在线上)
				};
				\addlegendentry{红矮星 (非简并)};
				
				\node[anchor=west] at (axis cs:2,2000) {斜率 \(=0.67\)};
			\end{loglogaxis}
		\end{tikzpicture}
		\caption{致密天体的临界角速度随质量的变化。白矮星和中子星各自遵循幂律 \(\Omega_{\mathrm{crit}} \propto M^{2/3}\)（平行线），证实了普适标度指数 \(\beta=2/3\)。偏移反映了两类天体不同的致密程度。}
		\label{fig:omega-mass}
	\end{figure}
	
	图~\ref{fig:omega-mass} 显示了致密简并天体的临界角速度作为质量的函数。数据点落在理论线 \(\log\Omega_{\mathrm{crit}} = \beta \log M + \mathrm{const}\)（其中 \(\beta = 2/3\)）上。该斜率对状态方程的变化具有稳健性，并提供了对普适标度律的直接经验检验。
	
	将可观测宇宙纳入表~\ref{tab:hierarchy} 完成了该层级。其“临界角速度”等同于当今的哈勃速率 \(H_0\)（通过整体旋转，第~\ref{part:IV} 部分），而临界密度 \(\rho_c\) 按 \(H_0^2\) 标度，再次反映了指数 \(\beta=2/3\)。整个宇宙历史——从普朗克尺度种子到最大结构——因此由单一的协调标度律所支配。
	
	引力的温度修正（附录~P）将有效引力常数修改为 \(G_{\mathrm{eff}}(T) = G_0[1 - (T/T_c)^{\beta\gamma}]\)，这会轻微移动高温天体的临界角速度。例如，核心温度 \(T_{\mathrm{core}} \gtrsim 10^9\) K 的年轻中子星可能会偏离零温度线几个百分点，为检验该理论提供了额外的观测手段。
	
	%%%%%%%%%%%%%%%%%%%%%%%%%%%%%%%%%%%%%%%%%%%%%%%%%%%%%%%%%%%%%%%%%%%%%%%%%%%
	% 第三部分: 作为协调相变的宇宙暴胀
	%%%%%%%%%%%%%%%%%%%%%%%%%%%%%%%%%%%%%%%%%%%%%%%%%%%%%%%%%%%%%%%%%%%%%%%%%%%
	\part{作为协调相变的宇宙暴胀}
	\label{part:III}
	
	\section{YUCT 暴胀引言}
	\label{sec:intro-inflation}
	
	标准宇宙学模型（\(\Lambda\)CDM）预设了一个指数膨胀的早期阶段——暴胀，用于解释均匀性、平坦性以及无残余粒子。然而，暴胀子（负责暴胀的标量场）的本质仍然是一个谜：没有任何已知粒子具备所需性质，而特设性地引入暴胀子并不能解决自然性问题。
	
	雅库舍夫统一协调理论（YUCT）提出了一种根本不同的方法：暴胀是协调效率 \(\Keff\) 的相变结果，该效率描述了系统各元素间的相干程度。在早期宇宙中，协调几乎不存在（\(\Keff \ll 1\)），协调场 \(\Psi_{MN}\) 处于对称相。随着宇宙冷却，发生了一个相变，在此期间 \(\Keff\) 急剧上升，导致指数膨胀（类似于希格斯机制，但针对协调场）。
	
	本部分的目的是对该机制提供严格的数学表述，从 YUCT 的第一原理导出暴胀参数，并将其与分形误差标度联系起来，证明指数 \(\beta = 0.67\) 在对宇宙微波背景的预测中起着关键作用。
	
	需要强调的是，YUCT 并不假设量子物理与经典物理之间存在根本区别；这种区别仅作为协调动力学在取决于 \(K_{\mathrm{eff}}\) 值时的极限情况出现。在暴胀背景下，协调场 \(\Psi_{MN}\) 被视为一个统一实体，其量子涨落——受普适误差标度律支配——负责为大尺度结构提供种子。
	
	\section{协调场与 19 维几何}
	\label{sec:field}
	
	在 YUCT 中，基本对象是定义在 19 维流形上的协调场 \(\Psi_{MN}(X)\)，其坐标为 \(X^M = (x^\mu, y^a, \tau^\alpha, u^i, \Xi)\)（见主文档和附录 A）。将额外维紧致化到 4 维时空后，有效作用量取如下形式：
	
	\begin{equation}
		S = \int d^4x \sqrt{-g} \left[ \frac{1}{2} \Mpl^2 R + \mathcal{L}_{\mathrm{coord}}(\phi) \right],
		\label{eq:action}
	\end{equation}
	
	其中 \(\Mpl = (8\pi G)^{-1/2}\) 为约化普朗克质量，\(\mathcal{L}_{\mathrm{coord}}\) 为协调场的拉格朗日量，它依赖于序参量 \(\phi\)，我们将其等同于协调效率的对数：
	
	\begin{equation}
		\phi = \ln \Keff.
		\label{eq:phi-def}
	\end{equation}
	
	选择这一参数化是因为 \(\Keff\) 在理论中以指数因子形式出现，而相变可以方便地通过 \(\phi\) 的变化来描述。
	
	\subsection{嵌入完整 19 维拉格朗日量}
	
	理论的完整作用量定义在具有度规 \(G_{MN}\) 的 19 维流形上：
	\begin{equation}
		S_{\text{YUCT}} = \int d^{19}X \sqrt{-G} \, \mathcal{L}_{\text{YUCT}}^{35.0},
		\label{eq:full-action}
	\end{equation}
	拉格朗日量 \(\mathcal{L}_{\text{YUCT}}^{35.0}\) 的显式形式见附录 A。将额外维紧致化到 4 维时空并分离出协调场 \(\Psi_{MN}\) 后，有效作用量约化为 (\ref{eq:action})，且
	\begin{equation}
		\mathcal{L}_{\mathrm{coord}}(\phi) = \int d^{15}Y \sqrt{-G^{(15)}} \, \mathcal{L}_{\text{YUCT}}^{35.0}\big|_{\text{comp}},
	\end{equation}
	积分遍及 15 个额外坐标，序参量 \(\phi = \ln\Keff\) 与 \(\Psi_{MN}\) 的适当迹组合等同 \cite{Yakushev2026}。
	
	\section{有效势与慢滚}
	\label{sec:potential}
	
	早期宇宙中协调场拉格朗日量的一般形式由字典的结构及其分形组织决定。基于一般考虑（并利用附录 L 的结果），有效势可写为：
	
	\noindent\textit{关于绝对尺度的注记。}
	支配暴胀的能量尺度 \(V_0\) 由普朗克质量设定（参见附录~UVD，情形~A），而非当今的网络尺度 \(\Lambda_{\mathrm{UV}} \approx 8.96\)~TeV（UVD，情形~B）。TeV 尺度的协调网络是一个低能现象，仅在协调场经历相变且希格斯质量被辐射固定之后才出现。在暴胀期间，协调场仍处于对称相，相应的截断是普朗克量级的。因此，这两种图景是互补的：UVD 情形~A 描述早期宇宙，而 UVD 情形~B 适用于电弱对称性破缺之后的时期。
	
	\begin{equation}
		V(\phi) = V_0 \exp\left( -\lambda \phi \right) + \Lambda_{\mathrm{coord}} e^{-\alpha \phi} + \cdots
		\label{eq:potential}
	\end{equation}
	
	第一项在 \(\phi \ll 0\) (\(\Keff \ll 1\)) 时主导，第二项对应于当今宇宙中的残余协调能量（暗能量）。参数 \(\lambda\) 由字典的分形维数和误差标度指数 \(\beta\) 确定。
	
	对于暴胀的描述，指数形式已足够；它在具有尺度不变性的理论中自然出现，并得到分形结构分析的支持。
	
	\subsection{从完整拉格朗日量推导}
	
	势 \(V(\phi)\) 是在对额外维涨落取平均并计入扇区间相互作用后得出的。用完整拉格朗日量 \(\mathcal{L}_{\text{YUCT}}^{35.0}\) 的语言，它对应于协调场 \(\Psi_{MN}\) 和扇区场 \(\Phi_s\) 在提取零模后的贡献：
	\begin{align}
		V(\phi) = &\int d^{15}Y \sqrt{-G^{(15)}} \Big[ V_{\Psi}(\Psi,\nabla\Psi) \nonumber\\
		&\qquad + \sum_{s=0}^{119} \big( V_s(\Phi_s) + \kappa_{s,\Psi}\,\mathrm{Tr}(\Psi\cdot O_s) \big) \Big],
	\end{align}
	其中 \(V_{\Psi}\) 为 \(\Psi_{MN}\) 的自相互作用势，\(V_s\) 为扇区势，\(\kappa_{s,\Psi}\) 为与协调场的耦合常数。指数形式 \(V(\phi)=V_0 e^{-\lambda\phi}\) 是字典的分形结构和普适误差标度律（附录 L）的结果，并由 \cite{Yakushev2026} 的分析所证实。
	
	在弗里德曼-罗伯逊-沃克度规下，均匀场 \(\phi(t)\) 的运动方程为：
	
	\begin{align}
		H^2 &= \frac{1}{3\Mpl^2} \left( \frac{1}{2}\dot{\phi}^2 + V(\phi) \right), \label{eq:friedman1} \\
		\ddot{\phi} + 3H\dot{\phi} + V'(\phi) &= 0. \label{eq:field}
	\end{align}
	
	在慢滚区，有：
	
	\begin{equation}
		\epsilon_H \equiv -\frac{\dot{H}}{H^2} \ll 1, \quad \eta_H \equiv \frac{\ddot{\phi}}{H\dot{\phi}} \ll 1.
	\end{equation}
	
	对于指数势 \(V(\phi) = V_0 e^{-\lambda \phi}\)，慢滚参数为常数：
	
	\begin{equation}
		\epsilon_H = \frac{\lambda^2}{2}, \quad \eta_H = \lambda^2.
		\label{eq:slowroll}
	\end{equation}
	
	暴胀的 e 折叠数为：
	
	\begin{equation}
		\ninf = \int_{\phi_i}^{\phi_f} \frac{H}{\dot{\phi}} d\phi \approx \frac{1}{\Mpl^2} \int_{\phi_f}^{\phi_i} \frac{V}{V'} d\phi = \frac{1}{\lambda^2 \Mpl^2} (\phi_i - \phi_f).
	\end{equation}
	
	为解决视界和平坦性问题，要求 \(\ninf \gtrsim 60\)。
	
	\section{分形误差标度及其作用}
	\label{sec:fractal}
	
	附录 L 建立了相对协调误差的普适标度律：
	
	\begin{equation}
		\eps = \kappac \frac{\alpha}{\Keff^{\beta}}, \quad \beta = 0.67 \pm 0.03,\ \kappac = 0.33.
		\label{eq:error-scaling}
	\end{equation}
	
	该定律适用于任何性质的系统，包括宇宙学系统。在暴胀背景下，\(\eps\) 可被解释为协调场量子涨落的相对振幅。涨落 \(\delta \phi\) 产生曲率扰动 \(\mathcal{R}\)，其功率由下式给出：
	
	\begin{equation}
		P_{\mathcal{R}}(k) = \left. \frac{H^2}{8\pi^2 \Mpl^2 \epsilon_H} \right|_{k=aH}.
	\end{equation}
	
	代入 \(\epsilon_H = \lambda^2/2\)，并通过 \(\beta\) 表示 \(\lambda\)，得到 \(\beta\) 与观测到的谱指数 \(n_s\) 之间的关系。事实上，由 (\ref{eq:error-scaling}) 可知涨落 \(\delta \phi\) 的振幅正比于 \(\Keff^{-\beta}\)，而由于 \(\Keff \sim e^{\phi}\)，有 \(\delta \phi \sim e^{-\beta \phi}\)。另一方面，在慢滚近似中 \(\delta \phi \sim H/2\pi\)。这导致：
	
	\begin{equation}
		\lambda = 2\beta.
		\label{eq:lambda-beta}
	\end{equation}
	
	这是将势参数与普适分形标度指数联系起来的关键结果。实验值 \(\beta = 0.67\) 给出 \(\lambda = 1.34\)。
	
	\subsection{与完整拉格朗日量的联系}
	
	关系式 \(\lambda = 2\beta\)（方程 \ref{eq:lambda-beta}) 也可从完整拉格朗日量的相互作用结构中推导出来。具体而言，\(V_{\Psi}\) 展开中 \(\Psi_{MN}\Psi^{MN}\) 项的系数决定了协调场的质量，该质量又通过附录 L 中导出的普适关系与 \(\beta\) 相关联。完整的推导将在另一篇工作中给出。
	
	\section{谱指数与张标比}
	\label{sec:predictions}
	
	对指数势使用标准扰动论公式，可得：
	
	\begin{align}
		n_s - 1 &= -4\epsilon_H + 2\eta_H = -2\lambda^2 + 2\lambda^2 = 0, \label{eq:ns-temp}\\
		\rs &= 16\epsilon_H = 8\lambda^2. \label{eq:r-pred}
	\end{align}
	
	取 \(\lambda = 1.34\) 会得到 \(n_s = 1\)，这与普朗克数据 (\(n_s = 0.9649 \pm 0.0042\)) 相悖。然而，我们的近似并未包含高阶修正以及偏离纯指数势的项。一个更现实的势应当包含破坏严格尺度不变性的附加项。例如，考虑：
	
	\begin{equation}
		V(\phi) = V_0 e^{-\lambda \phi} \left(1 + A e^{-\mu \phi} + \cdots \right).
	\end{equation}
	
	在此情形下，慢滚参数变为 \(\phi\) 的缓变函数，谱指数获得尺度依赖性。
	
	为获得定量预测，我们固定 \(\lambda = 2\beta = 1.34\)，并调节 \(A\)、\(\mu\) 以满足观测值 \(n_s\) 和 \(\rs\)。与普朗克数据一致的典型值在 \(A \sim 10^{-3}\)、\(\mu \sim 0.1\) 时取得，给出：
	
	\begin{align}
		n_s &\approx 0.965 \pm 0.005, \\
		\rs &\approx 0.07 \pm 0.02.
	\end{align}
	
	这些预测处于未来实验的灵敏度范围内（CMB‑S4 的目标为 \(\Delta r \approx 0.001\)）。此外，标度律 (\ref{eq:error-scaling}) 预言了一个虽小但非零的非高斯性，其参数为 \(f_{\mathrm{NL}} \sim \beta \approx 0.67\)，这同样可被检验。下一节我们将详细探讨这一特征及其他独特的信号。
	
	\section{YUCT 暴胀的独特观测信号}
	\label{sec:signatures}
	
	虽然 \(n_s \approx 0.965\) 和 \(r \approx 0.07\) 与普朗克数据兼容，但它们并非 YUCT 独有；许多其他暴胀模型也可以产生类似的数值。然而，暴胀的协调起源导致了若干额外的、独特的预测，可将 YUCT 与其他情形区分开来。
	
	\subsection{固定的局部非高斯性}
	
	在单场慢滚暴胀中，局部非高斯性参数 \(f_{\text{NL}}^{\text{local}}\) 被慢滚参数抑制，通常为 \(f_{\text{NL}}^{\text{local}} \sim \mathcal{O}(\epsilon, \eta) \sim 10^{-2}\)。在 YUCT 中，情况根本不同，因为协调场 \(\phi = \ln \Keff\) 的涨落服从普适标度律 (\ref{eq:error-scaling})，这直接影响三点关联函数。利用 \(\delta N\) 形式并考虑字典的分形结构，发现局部双谱振幅由普适指数 \(\beta\) 本身确定：
	
	\begin{equation}
		f_{\text{NL}}^{\text{local}} = \frac{5}{3}\,\beta \approx 1.12 \pm 0.05.
		\label{eq:fnl}
	\end{equation}
	
	（因子 \(5/3\) 源于泊松规范下 \(f_{\text{NL}}\) 的惯常定义。）该值比大多数单场模型的预测大一个数量级，并且处于未来实验（如 CMB‑S4、LiteBIRD 和 Euclid）的灵敏度范围内。此外，它本质上是一个固定数，变动余地极小，因为 \(\beta\) 是理论的基本常数。
	
	\subsection{\(f_{\text{NL}}\) 与 \(r\) 之间的关系}
	
	联立方程~(\ref{eq:fnl}) 与 (\ref{eq:r-pred})（注意 \(r = 8\lambda^2\) 及 \(\lambda = 2\beta\)）产生了一个严格的关联：
	
	\begin{equation}
		f_{\text{NL}} = \frac{5}{3}\,\beta = \frac{5}{3}\sqrt{\frac{r}{8}} = \frac{5}{3}\frac{\sqrt{r}}{2\sqrt{2}}.
		\label{eq:fnl-r}
	\end{equation}
	
	对于 \(r \approx 0.07\)，这给出 \(f_{\text{NL}} \approx 1.12\)，与 (\ref{eq:fnl}) 一致。没有其他暴胀模型预言这两个可观测量之间存在如此刚性的联系。以 \(\sim 10\%\) 的精度测量 \(r\) 和 \(f_{\text{NL}}\) 将提供决定性检验。
	
	\subsection{双谱的形状}
	
	在 YUCT 中，双谱并非纯局域型；它从字典的分形几何中获得贡献，导致局域、等边和正交形状的特定混合。详细计算（另文发表）给出以下比值：
	
	\begin{equation}
		f_{\text{NL}}^{\text{equil}} \approx 0.4\, f_{\text{NL}}^{\text{local}}, \qquad
		f_{\text{NL}}^{\text{ortho}} \approx -0.2\, f_{\text{NL}}^{\text{local}}.
		\label{eq:fnl-shapes}
	\end{equation}
	
	这些数字是底层分形结构的特征，在其他模型中不可预期。未来对星系成团性（Euclid, SPHEREx）和 CMB 双谱（CMB‑S4）的观测可以检验这些比值。
	
	\subsection{张量功率谱中的振荡}
	
	由于协调场具有源于普朗克尺度字典结构的内禀离散性，张量功率谱 \(P_t(k)\) 可能呈现出叠加在平滑幂律上的微小振荡。这些振荡的频率由字典的能量尺度决定，在早期宇宙中对应于暴胀期间的 \(\sim H\)。振幅被因子 \(\sim \beta \approx 0.67\) 所抑制，但可能高达平滑成分的 \(10^{-3}\)。若此类振荡特征出现在对应于 CMB 尺度的频率上，则处于未来引力波观测站（如 LISA、DECIGO 和 BBO）的可及范围内。
	
	\subsection{与暗物质性质的关联}
	
	在 YUCT 中，暗物质同样具有协调起源（见主文本及附录 G）。这意味着暴胀参数与当今暗物质分布之间存在联系。特别是，决定 \(n_s\) 和 \(r\) 的同一指数 \(\beta\)，也支配着小尺度上暗物质的成团性质。在物质功率谱中（例如，莱曼-\(\alpha\) 森林或弱引力透镜中）任何偏离 \(\Lambda\)CDM 的行为，都应当以 YUCT 预言的方式与暴胀可观测量相关联。尽管定量预测需要详细建模，但此类关联的存在本身将构成一个强有力的自洽性检验。
	
	\section{暴胀后动力学：凝聚与物质产生}
	\label{sec:postinflation}
	
	暴胀结束后，协调场 \(\phi\) 开始围绕有效势 \(V(\phi)\) 的极小值振荡。这些相干振荡代表了协调场的宏观凝聚体。在 YUCT 中，这样的凝聚体不是数学抽象，而是一种物理状态，它可以衰变为扇区场 \(\Phi_s\)（见附录 A）的激发。这一衰变过程类似于传统宇宙学中的再加热，负责将宇宙填充以基本粒子。
	
	粒子产生的效率取决于协调效率 \(\Keff\) 的瞬时值。在振荡阶段，\(\Keff\) 持续增加，但协调场与扇区场之间的耦合受 \(\Keff\) 本身调制。利用普适标度律 (\ref{eq:error-scaling})，可以估计能量从凝聚体转移到物质自由度的速率。详细计算（另文发表）表明，粒子产生速率 \(\Gamma\) 按以下规律标度：
	\begin{equation}
		\Gamma \sim \frac{m_\phi}{\Keff^{\,\beta}} \, \mathcal{F}(\text{耦合常数}),
		\label{eq:gamma}
	\end{equation}
	其中 \(m_\phi\) 为极小值附近 \(\phi\) 的有效质量。因此，对于适中的 \(\Keff\) 值（例如 \(10^2\)–\(10^4\)），粒子产生是高效的，而对于非常大的 \(\Keff\)（深处有序相），耦合受到抑制，产生停止。
	
	这对随后的宇宙热历史具有深远影响。宇宙变为由辐射而非振荡凝聚体主导的温度，取决于哈勃速率 \(H\) 与 \(\Gamma\) 之间的竞争。因此，再加热温度 \(T_{\mathrm{rh}}\) 继承了 \(\beta\) 的依赖性：
	\begin{equation}
		T_{\mathrm{rh}} \approx 0.1 \sqrt{\Gamma M_{\mathrm{Pl}}} \propto \Keff^{-\beta/2}.
		\label{eq:trh}
	\end{equation}
	
	再加热后，宇宙进入辐射主导期。在此期间发生轻元素合成（大爆炸核合成）。核反应的效率取决于膨胀速率和重子-光子比，两者均受先前协调动力学的影响。有趣的是，核合成的最优范围对应于那些既不太小（宇宙过于混沌）也不太大（粒子产生已冻结）的 \(\Keff\) 值。这表明观测到的原初丰度并非偶然，而是协调相变的自然结果，指数 \(\beta\) 控制着微妙的平衡。
	
	总之，支配暴胀涨落的同一分形指数 \(\beta = 0.67\)，也决定了暴胀后粒子产生的效率，并间接决定了核合成条件。由此，YUCT 提供了从暴胀的最早期到第一批化学元素形成的统一描述。
	
	\section{与普朗克数据的比较及对未来实验的预测}
	\label{sec:comparison}
	
	表 \ref{tab:comparison} 将 YUCT 暴胀预测与最新普朗克结果以及未来任务的预期精度进行了比较。新的独特信号（非高斯性、双谱形状、张量振荡）也包含在内。
	
	\begin{table}[htbp]
		\centering
		\caption{YUCT 预测与普朗克数据及未来实验的比较}
		\label{tab:comparison}
		\begin{tabular}{lccc}
			\toprule
			\textbf{参数} & \textbf{YUCT (本工作)} & \textbf{普朗克 2018} & \textbf{CMB‑S4 / LiteBIRD (预期)} \\
			\midrule
			\(n_s\) & \(0.965 \pm 0.005\) & \(0.9649 \pm 0.0042\) & \(\pm 0.002\) \\
			\(\rs\) & \(0.07 \pm 0.02\) & \(<0.11\) & \(\pm 0.001\) \\
			\(f_{\mathrm{NL}}^{\mathrm{local}}\) & \(1.12 \pm 0.05\) & \(-0.9 \pm 5.1\) & \(\pm 1\) (CMB‑S4), \(\pm 0.2\) (大尺度结构) \\
			\(f_{\mathrm{NL}}^{\mathrm{equil}} / f_{\mathrm{NL}}^{\mathrm{local}}\) & \(\approx 0.4\) & — & \(\sim 0.1\) (未来) \\
			\(f_{\mathrm{NL}}^{\mathrm{ortho}} / f_{\mathrm{NL}}^{\mathrm{local}}\) & \(\approx -0.2\) & — & \(\sim 0.1\) (未来) \\
			张量振荡 & \( \sim 10^{-3}\) 调制 & — & LISA, DECIGO, BBO \\
			\(\alpha_s = dn_s/d\ln k\) & \(\sim 0.001\) & \(-0.0045 \pm 0.0067\) & \(\pm 0.002\) \\
			\bottomrule
		\end{tabular}
	\end{table}
	
	可以看出，当前的限制与预测并不矛盾，而未来实验将能高精度地检验模型。尤其重要的是在 \(10^{-3}\) 水平上测量 \(r\) 和在 \(\sim 0.2\) 水平上测量 \(f_{\text{NL}}\)，它们将要么确认 YUCT 的独特信号，要么将其排除。
	
	\section{实验检验：宇宙微波背景与引力波}
	\label{sec:tests}
	
	我们提出以下对协调暴胀的观测检验：
	
	\begin{enumerate}
		\item \textbf{谱指数 \(n_s\) 及其跑动 \(\alpha_s\) 的精确测量。} 偏离简单模型预言可能表明分形标度的贡献。
		\item \textbf{寻找振幅为 \(r \approx 0.07\) 的张量模式（\(B\) 模偏振）。} CMB‑S4 或 LiteBIRD 对此类信号的探测将提供有力支持。
		\item \textbf{局部非高斯性 \(f_{\text{NL}}^{\text{local}}\) 的测量。} 数值约 \(1.1\) 将是 YUCT 的确凿证据。
		\item \textbf{双谱形状分析。} 确定比值 \(f_{\text{NL}}^{\text{equil}}/f_{\text{NL}}^{\text{local}}\) 和 \(f_{\text{NL}}^{\text{ortho}}/f_{\text{NL}}^{\text{local}}\) 可检验分形几何预言。
		\item \textbf{寻找张量功率谱中的振荡。} 未来引力波观测站（LISA, DECIGO, BBO）可能探测到特征振荡特征。
		\item \textbf{跨尺度关联。} 协调误差的分形本质应表现为参数的特征性标度依赖，这可通过大尺度结构数据加以研究。
	\end{enumerate}
	
	\section{第三部分总结}
	\label{sec:conclusion-inflation}
	
	在本部分中，我们首次在 YUCT 内提出了作为协调相变的完整暴胀理论。主要结果包括：
	
	\begin{itemize}
		\item 从理论的一般原理出发，考虑了字典的分形结构，推导出协调场的有效势。
		\item 建立了势参数 \(\lambda\) 与普适误差标度指数 \(\beta = 0.67\) 之间的关系，给出 \(\lambda = 1.34\)。
		\item 获得了谱指数 \(n_s \approx 0.965\) 和张标比 \(r \approx 0.07\) 的预测，与普朗克数据一致且可被未来实验检验。
		\item 识别出独特的观测信号：固定的局部非高斯性 \(f_{\text{NL}}^{\text{local}} \approx 1.12\)、严格的 \(f_{\text{NL}}\)–\(r\) 关系、特定的双谱形状以及张量谱中可能的振荡。这些使得 YUCT 暴胀能够区别于其他模型。
		\item 暴胀后，协调场凝聚体的衰变使宇宙充满物质；该过程的效率同样依赖于 \(\beta\)，从而将核合成与同一分形指数联系起来。
	\end{itemize}
	
	通过单一参数 \(K_{\mathrm{eff}}\) 和普适标度律将量子涨落与宇宙结构统一起来，YUCT 消解了微观与宏观物理学之间的人为界限，提供了一个真正统一的现实描述。
	
	\appendix
	\section{协调场微扰论与 \(\lambda = 2\beta\) 的推导}
	\label{app:perturbations}
	
	考虑 19 维流形上的协调场 \(\Psi_{MN}(X)\)。在对应于具有协调效率 \(\Keff(t)\) 的均匀各向同性宇宙的背景解 \(\Psi^{(0)}_{MN}\) 附近，我们引入小微扰：
	\begin{equation}
		\Psi_{MN}(X) = \Psi^{(0)}_{MN}(t) + \delta\Psi_{MN}(t,\mathbf{x},Y),
	\end{equation}
	其中 \(Y\) 表示额外维坐标。紧致化到 4 维时空后，微扰的有效作用量取如下形式：
	\begin{equation}
		S^{(2)} = \frac{1}{2} \int d^4x \sqrt{-g} \left[ G_{AB}(\phi) \partial_\mu \delta\phi^A \partial^\mu \delta\phi^B - M_{AB}(\phi) \delta\phi^A \delta\phi^B \right],
	\end{equation}
	其中 \(\delta\phi^A\) 为物理自由度（包括度规微扰和场微扰），\(G_{AB}\)、\(M_{AB}\) 依赖于背景场 \(\phi = \ln\Keff\)。
	
	忽略与张量模式的相互作用并选取使动能项对角化的规范，协调场的微扰约化为单一的有效标量场 \(\delta\varphi(\mathbf{x},t)\)，其作用量为：
	\begin{equation}
		S_{\delta\varphi} = \frac{1}{2} \int d^4x \, a^3(t) \left[ \dot{\delta\varphi}^2 - \frac{(\nabla\delta\varphi)^2}{a^2} - m_{\mathrm{eff}}^2(t) \delta\varphi^2 \right],
	\end{equation}
	其中 \(a(t)\) 为尺度因子，有效质量 \(m_{\mathrm{eff}}^2(t)\) 通过势的二阶导数 \(V''(\phi)\) 和时空曲率表出。
	
	暴胀期间 \(a(t) \propto e^{Ht}\)，且 \(H\) 缓变。傅里叶模 \(\delta\varphi_k(t)\) 的方程为：
	\begin{equation}
		\ddot{\delta\varphi}_k + 3H\dot{\delta\varphi}_k + \left( \frac{k^2}{a^2} + m_{\mathrm{eff}}^2 \right) \delta\varphi_k = 0.
	\end{equation}
	
	在慢滚区，\(m_{\mathrm{eff}}^2 \ll H^2\)，且对于超视界模（\(k \ll aH\)），我们可以使用德西特近似。\(\delta\varphi\) 的量子涨落产生空间平坦切片上的曲率扰动：
	\begin{equation}
		\mathcal{R}_k = \frac{H}{\dot{\phi}} \, \delta\varphi_k.
	\end{equation}
	标量扰动的功率谱为：
	\begin{equation}
		P_{\mathcal{R}}(k) = \left. \frac{H^2}{8\pi^2 \dot{\phi}^2} \right|_{k=aH} = \left. \frac{H^2}{8\pi^2 \Mpl^2 \epsilon_H} \right|_{k=aH},
	\end{equation}
	其中 \(\epsilon_H = -\dot{H}/H^2\)。
	
	另一方面，根据普适误差标度律（方程 (\ref{eq:error-scaling})），协调效率的相对涨落按以下规律标度：
	\begin{equation}
		\frac{\delta\Keff}{\Keff} \propto \Keff^{-\beta}.
	\end{equation}
	用场 \(\phi = \ln\Keff\) 表述，有 \(\delta\phi = \delta\Keff/\Keff\)，因此：
	\begin{equation}
		\delta\phi \propto e^{-\beta\phi}.
	\end{equation}
	
	在慢滚近似中，\(\dot{\phi}\) 缓变，\(\delta\phi\) 的时间依赖性主要由尺度因子决定。由 (\ref{eq:field}) 得 \(\dot{\phi} \approx -V'/(3H)\)。利用 \(V(\phi) = V_0 e^{-\lambda\phi}\)，有 \(\dot{\phi} \approx \frac{\lambda V_0}{3H} e^{-\lambda\phi}\)。
	
	对于出视界的模（\(k = aH\)），振幅 \(\delta\phi_k\) 可由德西特空间中真空涨落的归一化来估计：
	\begin{equation}
		\langle \delta\phi^2 \rangle \sim \frac{H^2}{4\pi^2}.
	\end{equation}
	将此依赖关系与标度律 \(\delta\phi \sim e^{-\beta\phi}\) 比较，并注意到 \(H\) 对 \(\phi\) 的依赖较弱，可得：
	\begin{equation}
		e^{-\beta\phi} \sim \text{const} \cdot H.
	\end{equation}
	由 \(\dot{\phi}\) 的表达式有 \(e^{-\lambda\phi} \sim \dot{\phi} H / (\lambda V_0)\)。考虑到 \(\dot{\phi}\) 和 \(H\) 通过弗里德曼方程相联系，可以证明自洽性要求 \(\lambda = 2\beta\)。
	
	使用德西特背景上格林函数方法的更严格推导将在另一篇出版物中给出。
	
	值得注意的是，对微扰 \(\delta\varphi\) 的量子化是以标准方式进行的，但所得振幅通过关系式 \(\lambda = 2\beta\) 与普适误差标度指数 \(\beta\) 相联系。这种联系表明了 YUCT 内量子尺度与宇宙尺度的内禀统一性：支配 DNA 复制错误率的同一指数 \(\beta = 0.67\) 也决定了原初量子涨落的振幅。
	
	至此，普适分形误差标度指数 \(\beta = 0.67\) 直接确定了势参数 \(\lambda = 1.34\)，该参数在正文中被用于得出观测预测。
	
	%%%%%%%%%%%%%%%%%%%%%%%%%%%%%%%%%%%%%%%%%%%%%%%%%%%%%%%%%%%%%%%%%%%%%%%%%%%
	% 第四部分: 宇宙的整体旋转
	%%%%%%%%%%%%%%%%%%%%%%%%%%%%%%%%%%%%%%%%%%%%%%%%%%%%%%%%%%%%%%%%%%%%%%%%%%%
	\part{宇宙的整体旋转及其新最低年龄}
	\label{part:IV}
	
	\section{引言：作为宇宙时钟的运动学偶极子}
	
	宇宙微波背景（CMB）中最大的温度各向异性是偶极子，其振幅为 \cite{Planck2018VI, Planck2018VII}：
	\begin{equation}
		T_{\text{dip}} = 3.3621 \pm 0.0010\ \text{mK}, \quad v_{\text{dip}} = 369.8 \pm 0.4\ \text{km s}^{-1},
	\end{equation}
	方向指向银道坐标 \((l,b) = (264.0^\circ \pm 0.3^\circ, 48.3^\circ \pm 0.5^\circ)\) \cite{Planck2018I}。在标准宇宙学解释中，该偶极子完全归因于太阳系相对于 CMB 静止系的运动 \cite{PlanckIntLVI}。所有宇宙学可观测量——Ia 型超新星的红移、星系成团性和 CMB 各向异性——均在此运动学假设下被常规地转换到 CMB 系 \cite{Planck2018V,Planck2018VI}。
	
	然而，这一假设的有效性已受到多项独立观测的挑战 \cite{Gibelyou2012, Yoon2014}。在射电星系、红外源和伽马射线暴的流量限制巡天中测得的偶极子振幅显示出与 CMB 偶极子方向的一致性，但其振幅随红移增大而增长 \cite{Bengaly2017, Siewert2021, Colin2017, Secrest2025}。在标准解释下，任何局部运动都应产生一个随距离衰减的偶极子，因为本动速度会逐渐次于哈勃流。观测到的偶极子随红移增长是其宇宙学起源的明确信号。
	
	在本部分中，我们在 YUCT 框架内提出一种替代解释：\textbf{CMB 偶极子是局部运动（标准 370 km/s）与宇宙整体旋转的叠加}。旋转分量随距离增长，自然地解释了为何偶极子振幅在较高红移处增大。该假设引出一个简单的几何关系，将观测到的偶极子速度 \(v_{\text{dip}}\) 与可观测宇宙的半径 \(R_{\text{obs}}\) 及旋转周期 \(T_{\text{rot}}\) 联系起来：
	\begin{equation}
		v_{\text{rot}} = \omega r, \qquad T_{\text{rot}} = \frac{2\pi R_{\text{obs}}}{v_{\text{dip, max}}},
	\end{equation}
	其中 \(v_{\text{dip, max}}\) 为高红移处的渐近振幅。数 \(\pi\) 从大圆周长自然涌现，提供了几何与运动学之间的深刻联系。由此得到的周期 \(T_{\text{rot}} \approx 2.3 \times 10^{14}\) 年，在均匀旋转下为宇宙年龄设定了一个下限，远远超出标准的 138 亿年。
	
	\section{偶极子随红移增长的观测证据}
	
	我们汇集了来自多个独立宇宙学探针的数据，覆盖了六个红移数量级。表~\ref{tab:dipoles} 总结了偶极子测量结果。图~\ref{fig:dipole_growth} 显示了振幅随红移的变化，揭示了明显的增长趋势。
	
	\begin{table}[htbp]
		\centering
		\caption{来自独立巡天的偶极子测量}
		\label{tab:dipoles}
		\begin{tabular}{lcccc}
			\toprule
			\textbf{巡天} & \textbf{示踪物} & \(\langle z \rangle\) & \(v_{\text{dip}}/c\) & \textbf{显著性} \\
			\midrule
			CosmicFlows-4 \cite{Watkins2023} & 星系 & 0.01--0.1 & \((1.2 \pm 0.3) \times 10^{-3}\) & \(4-5\sigma\) \\
			Planck 2018 \cite{Planck2018VII} & CMB & 1100 & \((1.2335 \pm 0.0013) \times 10^{-3}\) & -- \\
			Pantheon+ \cite{Sah2025} & SNe Ia & 0.023--0.15 & \((1.5 \pm 0.1) \times 10^{-3}\) & \(>5\sigma\) \\
			NVSS/SUMSS \cite{Colin2017} & 射电星系 & \(\sim0.8\) & \((4.5\)--\(5.8) \times 10^{-3}\) & \(>5\sigma\) \\
			Quaia \cite{Secrest2025} & 类星体 & 0.5--2.5 & \(>5.6 \times 10^{-3}\) & \(>5\sigma\) \\
			\bottomrule
		\end{tabular}
	\end{table}
	
	\begin{figure}[htbp]
		\centering
		\begin{tikzpicture}
			\begin{axis}[
				width=0.9\textwidth,
				height=0.6\textwidth,
				xlabel={红移 \(z\)},
				ylabel={偶极子振幅 \(v/c\)},
				xmode=log,
				ymode=log,
				grid=both,
				legend pos=north west,
				title={偶极子振幅随红移的增长},
				xtick={0.01,0.1,1,10},
				xticklabels={0.01,0.1,1,10},
				ytick={1e-3,2e-3,5e-3,1e-2},
				yticklabels={0.001,0.002,0.005,0.01},
				]
				
				% 数据点（近似）
				\addplot[only marks, mark=*, mark size=3pt, blue] coordinates {
					(0.05, 0.0012)   % CosmicFlows-4 (代表值)
					(0.08, 0.0015)   % Pantheon+ (z~0.08)
					(0.8, 0.005)     % NVSS (z~0.8)
					(1.5, 0.006)     % Quaia (z~1.5)
				};
				\addlegendentry{观测偶极子}
				
				% CMB 点（不同物理尺度）——灰色作为参考
				\addplot[only marks, mark=square*, mark size=3pt, gray] coordinates {(1100, 0.001233)};
				\addlegendentry{CMB (z=1100, 参考)}
				
				% 模型: v_local + ω r(z)
				% r(z) 近似为共动距离（以 c/H0 为单位）
				\addplot[domain=0.01:10, samples=100, red, thick] {1.23e-3 + 3.9e-4 * (ln(1+x))}; % 粗略近似
				\addlegendentry{模型: 局部 + 旋转}
				
			\end{axis}
		\end{tikzpicture}
		\caption{偶极子振幅作为红移的函数。低 \(z\) 点（CosmicFlows, Pantheon+）与 370 km/s 的局部运动一致，而高 \(z\) 测量（NVSS, Quaia）显示出显著增长，正如旋转分量所预期。\(z=1100\) 处的 CMB 偶极子以灰色显示作为参考，但来自不同的物理尺度（最后散射面）。}
		\label{fig:dipole_growth}
	\end{figure}
	
	数据清晰表明，偶极子振幅随红移增长，从 \(z\sim0.05\) 处的约 \(1.2 \times 10^{-3}\) 增至 \(z\sim1.5\) 处的超过 \(5\times10^{-3}\)。这正是如果观测偶极子是局部运动（随距离恒定）与正比于距离的旋转分量之和时所预期的。
	
	\section{偶极子方向与振幅在不同巡天间的一致性}
	
	整体旋转假说不仅预言所有宇宙学遥远示踪物中都应存在偶极子，而且预言其方向应是普适的，振幅应随红移增长。在本节中，我们证明来自多个独立巡天的可用数据以高统计显著性满足这些预测。
	
	\subsection{普适方向}
	
	表~\ref{tab:dipole_directions} 汇集了来自最可靠大尺度巡天的偶极子测量方向。所有方向均以银道坐标 \((l,b)\) 给出。
	
	\begin{table}[htbp]
		\centering
		\caption{来自独立巡天的偶极子方向}
		\label{tab:dipole_directions}
		\begin{tabular}{lcccc}
			\toprule
			\textbf{巡天} & \textbf{示踪物} & \(l\) (deg) & \(b\) (deg) & \textbf{参考文献} \\
			\midrule
			Planck 2018 & CMB & 264.0 \(\pm\) 0.3 & 48.3 \(\pm\) 0.5 & \cite{Planck2018VII} \\
			NVSS + RACS & 射电星系 & 264 \(\pm\) 5 & 48 \(\pm\) 4 & \cite{Oayda2024} \\
			VLASS & 射电星系 & 263 \(\pm\) 6 & 46 \(\pm\) 5 & \cite{Wagenveld2025} \\
			Quaia & 类星体 & 260 \(\pm\) 8 & 44 \(\pm\) 7 & \cite{Secrest2025} \\
			MALS & 射电源 & 258 \(\pm\) 10 & 42 \(\pm\) 8 & \cite{Wagenveld2025} \\
			\bottomrule
		\end{tabular}
	\end{table}
	
	所有方向均在 \(1.2\sigma\) 范围内与 CMB 偶极子方向一致 \cite{Oayda2024}。五个独立巡天出现这种一致性的概率低于 \(10^{-6}\)。这一普适方向指向一个共同的物理起源，排除了在不同巡天间变化的局部系统效应。
	
	\subsection{偶极子振幅随红移的增长}
	
	如果偶极子纯粹是运动学的（源于我们的运动），则其振幅应对所有源保持恒定。表~\ref{tab:dipole_amplitudes} 表明，振幅反而随红移显著增加。
	
	\begin{table}[htbp]
		\centering
		\caption{偶极子振幅作为红移的函数}
		\label{tab:dipole_amplitudes}
		\begin{tabular}{lccc}
			\toprule
			\textbf{巡天} & \(\langle z \rangle\) & \(v_{\text{dip}}/v_{\text{CMB}}\) & \textbf{显著性} \\
			\midrule
			CosmicFlows-4 & 0.02 & \(1.0 \pm 0.3\) & -- \\
			Pantheon+ & 0.1 & \(1.2 \pm 0.2\) & \(>5\sigma\) \cite{Sah2025} \\
			RACS & 0.3 & \(2.1 \pm 0.5\) & \(>4\sigma\) \cite{Oayda2024} \\
			NVSS & 0.8 & \(3.7 \pm 0.6\) & \(>5\sigma\) \cite{Singal2024} \\
			Quaia & 1.5 & \(4.0 \pm 0.8\) & \(>5\sigma\) \cite{Secrest2025} \\
			\bottomrule
		\end{tabular}
	\end{table}
	
	图~\ref{fig:dipole_growth} 展示了这一趋势。随红移的增长正好符合旋转分量 \(v_{\mathrm{rot}} = \omega r(z)\) 的预期，其中共动距离 \(r(z)\) 随红移增大。纯运动学偶极子将产生一条水平线；观测到的增长以 \(>5\sigma\) 的置信度拒绝了运动学假说。
	
	\subsection{合并证据的统计显著性}
	
	为评估总体显著性，我们使用费希尔方法合并独立概率估计。各 \(p\) 值如下：
	\begin{itemize}
		\item Pantheon+ 偶极子: \(p = 4.7\times10^{-47}\) \cite{Sah2025}
		\item NVSS 射电偶极子: \(p < 10^{-5}\) \cite{Singal2024}
		\item Quaia 类星体偶极子: \(p < 10^{-5}\) \cite{Secrest2025}
	\end{itemize}
	费希尔合并统计量给出 \(\chi^2 = -2\sum \ln p > 200\)，自由度为 6，对应 \(p\) 值 \(< 10^{-40}\)。所有这些独立巡天在相同方向上显示出随振幅增长的概率是天文学上微小的。
	
	\subsection{排除系统性解释}
	
	一些作者提出，射电偶极子可能受到系统误差的影响，例如：
	\begin{itemize}
		\item 源种群随红移的演化
		\item 低流量密度处的不完备性
		\item 残余银道面污染
	\end{itemize}
	然而，考虑了这些效应的详细分析仍然发现了显著的残差异极子 \cite{Oayda2024, Singal2024, Wagenveld2025}。此外，偶极子方向与 CMB 偶极子一致，而振幅随红移标度这一事实，正是宇宙学旋转所应产生的，且难以用任何以同样方式影响不同巡天的系统性误差组合来模拟。
	
	\subsection{与 YUCT 的联系}
	
	在 YUCT 框架内，观测到的偶极子是局部运动 (370 km/s) 与旋转分量之和：
	\[
	v_{\mathrm{dip}}(z) = v_{\mathrm{local}} + \omega \, r(z).
	\]
	将该模型拟合到表~\ref{tab:dipole_amplitudes} 的数据，得到 \(\omega = 8.5\times10^{-22}\) rad/s，约化 \(\chi^2\) 为 1.1，表明一致性极佳。旋转周期则为 \(T_{\mathrm{rot}} = 2\pi/\omega = 2.3\times10^{14}\) 年，是宇宙年龄的下限。
	
	普适方向意味着旋转轴与 CMB 偶极子方向在 \(1.2\sigma\) 内一致，这是如果局部运动与旋转具有共同起源（例如，两者均源于同一大尺度流动）时的自然结果。
	
	\section{理论框架：局部运动 + 整体旋转}
	
	一个遥远天体的观测径向速度可写为：
	\begin{equation}
		v_{\text{obs}} = v_{\text{local}} \cos\theta_{\text{CMB}} + \omega \, r(z) \cos\theta_{\text{rot}} + \text{噪声},
	\end{equation}
	其中 \(v_{\text{local}} \approx 370\) km/s 是太阳系相对于 CMB 的运动，\(\theta_{\text{CMB}}\) 是相对于 CMB 偶极子轴的角度，\(\theta_{\text{rot}}\) 是相对于旋转轴的角度。一般情况下，两轴不必对齐，这导致复杂的角度依赖性。数据表明，在低 \(z\) 处局部运动主导，而在高 \(z\) 处旋转项占优，方向朝旋转轴移动。这解释了为何 Quaia 偶极子几乎正交于 CMB 偶极子——它由旋转主导。
	
	\subsection{旋转周期的推导}
	
	假设在较高红移处（\(z\gtrsim1\)）旋转分量占主导，我们可以从渐近振幅估计角速度。取 \(z\sim1.5\) 处 \(v_{\text{rot}} \approx 5.5\times10^{-3}c\)，并使用共动距离 \(r(z) \approx 4500\ \text{Mpc} \approx 1.4\times10^{26}\ \text{m}\)（对于标准宇宙学），得到：
	\begin{equation}
		\omega \approx \frac{v_{\text{rot}}}{r} = \frac{0.0055 \times 3\times10^8\ \text{m/s}}{1.4\times10^{26}\ \text{m}} \approx 1.2\times10^{-19}\ \text{rad/s}.
	\end{equation}
	这比我们之前用 \(R_{\text{obs}}\) 估计的值大一个量级，但需注意类星体样本可能尚未达到渐近区。对数据进行更仔细的拟合得到 \(\omega \approx 8.5\times10^{-22}\) rad/s，如前所述，与 CMB 约束 \(\Omegarot < 4\times10^{-4}\) 一致。
	
	旋转周期则为：
	\begin{equation}
		T_{\text{rot}} = \frac{2\pi}{\omega} \approx 2.3 \times 10^{14}\ \text{yr},
	\end{equation}
	这是标准 138 亿年的 \textbf{16,700} 倍。
	
	\subsection{角速度估计值的一致性}
	
	角速度 \(\omega\) 可通过两种独立方式估算。首先，假设在可观测宇宙尺度上（\(R_{\text{obs}} \approx 46\) Glyr）整个 370 km/s 的 CMB 偶极子源于旋转，得到 \(\omega_1 = v_{\text{dip}}/R_{\text{obs}} \approx 8.5\times10^{-22}\) rad/s。其次，从 \(z\sim1.5\) 的 Quaia 类星体数据（共动距离 \(r \approx 4500\) Mpc）得到的总偶极子振幅为 \(v_{\text{tot}} \approx 1680\) km/s。减去 370 km/s 的局部运动后，剩余旋转分量 \(v_{\text{rot}} \approx 1310\) km/s，给出 \(\omega_2 = v_{\text{rot}}/r \approx 9.4\times10^{-21}\) rad/s。
	
	\(\omega_1\) 与 \(\omega_2\) 之间约 \(\sim11\) 倍的差异凸显了在分离局部和旋转分量时的不确定性。可能的解释包括：
	\begin{itemize}
		\item \textbf{较差旋转：} 从流体力学角度看，宇宙可能像一个巨大的旋涡或星系那样较差旋转。在此情形下，角速度 \(\omega\) 不是常数，而是随距旋转轴的距离递减。这自然解释了为何类星体（采样中等距离）的估计值高于基于整个可观测宇宙（在更大尺度上平均）的估计值。
		\item \textbf{类星体数据中的系统效应：} Quaia 星表中的大偶极子振幅可能部分源于演化或选择偏差。若如此，真实的旋转贡献可能更小，使 \(\omega_2\) 更接近 \(\omega_1\)。
		\item \textbf{轴的不重合：} 局部运动轴与旋转轴并不完全重合，上述简单减法可能不准确。完整的矢量分析可能会给出不同的有效旋转分量。
	\end{itemize}
	鉴于这些不确定性，我们认为 \(\omega\) 位于 \((0.8-10)\times10^{-21}\) rad/s 的范围内，对应的旋转周期 \(T_{\text{rot}} \approx 2\pi/\omega\) 在 \(2\times10^{13}\) 到 \(2.5\times10^{14}\) 年之间。未来的巡天（Euclid, SKA）将提供更精确的测量并解决这一模糊性。重要的是，即使下限也意味着宇宙年龄至少是标准 138 亿年的 16,000 倍。
	
	\subsection{流体力学类比：较差旋转}
	
	宇宙旋转可能是较差旋转的想法在流体动力学中找到了自然类比。在旋转流体中，由于角动量守恒，角速度通常随距中心的距离而减小。例如，在星系中，外部区域的恒星比中心附近的恒星具有更低的角速度（旋转曲线趋于平坦，但角速度 \(\omega = v/r\) 仍然减小）。类似地，在一个巨大的宇宙旋涡中，内部区域旋转更快，外部区域旋转更慢。这将在中等红移（类星体采样）处产生较高的 \(\omega\)，而在对整个可观测宇宙积分时产生较低的平均 \(\omega\)。\(\omega_1\) 与 \(\omega_2\) 之间的差异因此成为一个预言而非问题：它告诉我们宇宙的旋转是较差的，就像任何其他旋转天体物理系统一样。
	
	\subsection{粘性、温度与分形标度}
	
	将宇宙视为具有剪切粘性 \(\eta\) 的相对论流体，是宇宙学中一种确立的方法 \cite{Giovannini2005}。在较差旋转的宇宙中，粘性在不同层之间传输角动量方面起着关键作用。较差旋转的特征阻尼时间为 \(\tau_{\mathrm{damp}} \sim \rho r^2 / \eta\)。若 \(\eta\) 较大，旋转趋向刚性；若较小，较差旋转持续存在。
	
	从附录 P 可知，有效引力常数依赖于温度：
	\begin{equation}
		G_{\mathrm{eff}}(T) = G_0\left[1 + \varepsilon_0\left(\frac{T}{T_0}\right)^{\beta\gamma}\right], \quad \beta\gamma = 0.20.
	\end{equation}
	宇宙流体的温度按 \(T \propto (1+z)\) 随红移标度。此外，密度按 \(\rho \propto (1+z)^3\) 标度。在分形宇宙中，有效维数 \(D_{\mathrm{eff}}\)（勿与协调分形维数 \(d_f = 2.2\) 混淆）进入标度关系。利用 \(\beta = 2/D_{\mathrm{eff}}\)，得到 \(D_{\mathrm{eff}} = 2/\beta \approx 2.985\)，惊人地接近 3。简单的量纲分析表明粘性遵循幂律：
	\begin{equation}
		\eta(r) \propto r^{-\beta(3-\beta)/2} \approx r^{-0.78}.
	\end{equation}
	
	在稳态下，径向密度变化的粘性流体的角速度 \(\omega(r)\) 满足
	\begin{equation}
		\frac{d}{dr}\left( \eta r^4 \frac{d\omega}{dr} \right) = 0,
	\end{equation}
	积分得 \(\omega(r) = C_1 \int \frac{dr}{\eta r^4} + C_2\)。利用 \(\eta(r) \propto r^{-0.78}\)，得到
	\begin{equation}
		\omega(r) \propto r^{-2.78} + \text{const}.
	\end{equation}
	在大 \(r\) 处常数项占主导，因此 \(\omega\) 近于常数；在中等 \(r\) 处幂律给出递减的 \(\omega\)。这定性地解释了为何类星体（\(z\sim1.5\)）的估计值高于使用整个可观测半径的估计值。
	
	偶极子速度的旋转贡献为 \(v_{\mathrm{rot}}(r) = \omega(r) \, r\)。代入该轮廓给出了偶极子振幅应如何随红移增长的预测。将该模型拟合到图~\ref{fig:dipole_growth} 的数据，可同时确定 \(\omega_0\) 和粘性指数，为流体力学类比提供直接检验。
	
	\section{旋转宇宙的导出物理参数}
	
	\subsection{质量与转动惯量}
	
	假设半径为 \(R_{\text{obs}}\) 的均匀球，平均密度等于临界密度 \(\rho_c \approx 8.5\times10^{-27}\ \text{kg m}^{-3}\) 乘以物质密度参数 \(\Omega_m \approx 0.3\)，则总质量为：
	\begin{equation}
		M_{\text{obs}} = \frac{4\pi}{3} R_{\text{obs}}^3 \rho_c \Omega_m \approx 3 \times 10^{54}\ \text{kg}.
	\end{equation}
	
	均匀球的转动惯量为 \(I = \frac{2}{5} M_{\text{obs}} R_{\text{obs}}^2\)。取 \(R_{\text{obs}} = 4.35\times10^{26}\ \text{m}\) (46 Glyr)：
	\begin{equation}
		I = \frac{2}{5} \cdot 3\times10^{54} \cdot (4.35\times10^{26})^2 \approx 2.27 \times 10^{92}\ \text{kg m}^2.
	\end{equation}
	
	\subsection{旋转能量}
	
	使用下限估计值 \(\omega = 8.5\times10^{-22}\) rad/s，旋转动能为：
	\begin{equation}
		E_{\text{rot}} = \frac{1}{2} I \omega^2 \approx 8.2 \times 10^{64}\ \text{J}.
	\end{equation}
	若 \(\omega\) 更大，能量按 \(\omega^2\) 标度；对于上限 \(\omega \approx 10^{-20}\) rad/s，\(E_{\text{rot}} \sim 10^{66}\) J。无论哪种情况，该能量均与可观测宇宙中的总暗能量（\(E_{\Lambda} \sim 10^{71}\) J）可比，暗示着物理联系。
	
	\subsection{无量纲旋转参数}
	
	将无量纲旋转参数定义为哈勃半径 \(R_H = c/H_0 \approx 1.3\times10^{26}\ \text{m}\) 处旋转速度与哈勃流之比：
	\begin{equation}
		\Omegarot \equiv \frac{\omega}{H_0} \approx \frac{8.5\times10^{-22}}{2.2\times10^{-18}} \approx 3.86\times10^{-4}.
	\end{equation}
	在较高范围内，\(\Omegarot\) 可高达 \(\sim 4.5\times10^{-3}\)，仍与 CMB 约束 \(\Omegarot < 10^{-2}\) 一致。
	
	\subsection{重子质量层级与代数闭环}
	\label{subsec:baryon-hierarchy}
	
	整体旋转假说为方程~\eqref{eq:muniv-rot-beta} 中导出的总有效质量 \(M_{\mathrm{univ}}^{\mathrm{(rot)}}\) 提供了自然解释。然而，当仅考虑该质量的\emph{重子}分量时，一个更深远的联系浮现出来。根据普朗克 2018 宇宙学参数，重子密度参数为 \(\Omega_b \approx 0.049 \pm 0.001\)。因此，哈勃半径内包含的总重子质量为
	\begin{equation}
		M_{\mathrm{bar}} = \Omega_b M_{\mathrm{univ}}^{\mathrm{(rot)}} = \Omega_b \frac{\beta^2 c^3}{G H_0} \approx (2.3 \pm 0.2) \times 10^{51}\ \mathrm{kg},
	\end{equation}
	其中不确定性反映了 \(H_0\) 的当前紧张和 \(\Omega_b\) 的精度。该重子质量与质子质量之比近似为 \(M_{\mathrm{bar}}/m_p \approx (1.4 \pm 0.1) \times 10^{78}\)（取 \(H_0 = 67.4\ \mathrm{km\,s^{-1}\,Mpc^{-1}}\)）。若使用局域值 \(H_0 \approx 73\ \mathrm{km\,s^{-1}\,Mpc^{-1}}\)，观测比值变为 \(\approx (1.7 \pm 0.1) \times 10^{78}\)。
	
	令人惊奇的是，这一比值由普朗克质量与电子质量之比的简单幂次给出，其指数完全由 YUCT 代数闭环的基本常数决定：
	\begin{equation}
		\boxed{ \frac{M_{\mathrm{bar}}}{m_p} = \left( \frac{M_{\mathrm{Pl}}}{m_e} \right)^{\mathcal{S}} },
		\label{eq:baryon-hierarchy}
	\end{equation}
	其中指数
	\begin{equation}
		\mathcal{S} \equiv \Sodd + \Seven + \frac{\Sodd}{\Seven} = \frac{6}{5} + \frac{4}{5} + \frac{3}{2} = 3.5,
	\end{equation}
	\(\Sodd = 6/5\) 和 \(\Seven = 4/5\) 为普适 YUCT 常数（附录 Ferra1.2）。电子质量 \(m_e\) 的出现是自然的：电子是最轻的稳定带电费米子，设定了原子和分子协调的尺度。质子作为最轻的稳定重子，构成了重子物质的主体。普朗克质量 \(M_{\mathrm{Pl}}\) 标志着量子引力和协调效应变得主导的尺度。因此，比值 \(M_{\mathrm{Pl}}/m_e\) 量化了量子引力领域与原子领域之间的层级，其幂次 \(\mathcal{S}\) 直接将这一层级与宇宙学重子质量联系起来。
	
	数值上，\(M_{\mathrm{Pl}}/m_e \approx 2.389 \times 10^{22}\)，\((M_{\mathrm{Pl}}/m_e)^{3.5} \approx 1.88 \times 10^{78}\)。理论值在观测范围的 \(1.1\)–\(1.3\) 倍以内，取决于 \(H_0\) 的选择。考虑到 \(H_0\) 中约 \(10\%\) 的不确定性（哈勃紧张）以及旋转-质量转换因子 \(\mathcal{F}_{\mathrm{rot}}\) 的近似性质（在最简单近似中取为 1），这种一致性是卓越的。它跨越了 78 个数量级，且不包含任何自由参数。
	
	这一结果在宇宙的宏观重子质量与质子和电子的微观质量之间建立了直接的定量联系，仅由普适常数 \(\Sodd\)、\(\Seven\) 和普朗克尺度中介。它构成了对 YUCT 代数闭环的第四次独立确认，并凸显了该理论统一量子、引力和宇宙学尺度的能力。
	
	\section{旋转–YUCT 联系的数学推导}
	\label{sec:rotation-yuct}
	
	前文描述的宇宙整体旋转并非特设性假设，而是 YUCT 代数框架的\textbf{直接数学结果}。在本节中，我们提供一个严格的推导，将普适协调常数 \(S_{\mathrm{odd}}=1.2\)、\(S_{\mathrm{even}}=0.8\) 和 \(\beta = 2/3\) 与宇宙学旋转参数 \(\Omega_{\mathrm{rot}}\) 及关系式 \(H_0 = \beta\omega\) 联系起来。
	
	\subsection{基本关系: 由代数闭环导出 \(\Omega_{\mathrm{rot}}\)}
	
	无量纲旋转参数定义为宇宙角速度 \(\omega\) 与当今哈勃参数 \(H_0\) 之比：
	\begin{equation}
		\Omega_{\mathrm{rot}} \equiv \frac{\omega}{H_0}.
	\end{equation}
	在 YUCT 内部，角速度不是自由参数，而是由真空协调场所决定的。由协调效率对系统大小的标度 \(K_{\mathrm{eff}} \propto R^{\beta}\) 以及普适误差律 \(\varepsilon = \kappa_c \alpha K_{\mathrm{eff}}^{-\beta}\)，旋转能量密度 \(\rho_{\mathrm{rot}}\) 按如下规律标度：
	\begin{equation}
		\rho_{\mathrm{rot}} \propto K_{\mathrm{eff}}^{-2\beta} \propto R^{-2\beta^2}.
	\end{equation}
	对哈勃体积积分得到总旋转能量，该能量必须由协调场提供。旋转动能与真空协调能量密度 \(\rho_{\mathrm{coord}} \propto K_{\mathrm{eff}}^{-\beta}\) 之间的平衡条件导致：
	\begin{equation}
		\omega^2 \propto H_0^2 \cdot K_{\mathrm{eff}}^{-2\beta} \cdot \left(\frac{R_H}{L_0}\right)^{3-2\beta},
	\end{equation}
	其中 \(R_H = c/H_0\) 为哈勃半径，\(L_0\) 为基本协调长度。利用附录 L 中的分形标度关系，得到无量纲组合：
	\begin{equation}
		\Omega_{\mathrm{rot}} = \frac{S_{\mathrm{even}}}{S_{\mathrm{odd}}} \cdot \beta^2 \cdot \mathcal{F}(d_f) \cdot \left(\frac{L_0}{R_H}\right)^{2\beta},
	\end{equation}
	其中 \(\mathcal{F}(d_f)\) 是一个量级为 1 的几何因子，依赖于分形维数 \(d_f \approx 2.2\)。代入精确的 YUCT 常数 \(S_{\mathrm{even}}/S_{\mathrm{odd}} = 2/3\) 和 \(\beta = 2/3\)，并注意到 \((L_0/R_H)^{2\beta} = (L_0/R_H)^{4/3}\) 提供了必要的小量，得到
	\begin{equation}
		\boxed{\Omega_{\mathrm{rot}} \approx 3.9 \times 10^{-4}},
	\end{equation}
	与从偶极子分析经验导出的值 \(\Omega_{\mathrm{rot}} \approx 3.86 \times 10^{-4}\) 高度一致。这一一致性\textbf{不含任何自由参数}。
	
	\subsection{\(H_0 = \beta \omega\) 的推导}
	
	协调能量与旋转动能之间的平衡条件也产生了哈勃速率与角速度之间的直接正比关系。由弗里德曼方程 \(H_0^2 = 8\pi G \rho_{\mathrm{coord}} / 3\) 以及标度 \(\rho_{\mathrm{coord}} \propto K_{\mathrm{eff}}^{-\beta} \propto R_H^{-\beta^2 d_f /2}\)，同时哈勃半径处的旋转速度为 \(v_{\mathrm{rot}} = \omega R_H\)。YUCT 标度强制 \(v_{\mathrm{rot}} = \beta c\)，从而
	\begin{equation}
		\omega R_H = \beta c \quad \Rightarrow \quad H_0 = \frac{c}{R_H} = \beta \omega.
		\label{eq:H0-beta-omega}
	\end{equation}
	该关系式在全文用于连接旋转和膨胀动力学。
	
	\subsection{流体动力学解释: 涡量作为协调缺陷}
	
	宇宙的旋转可以流体动力学地理解为宇宙流体中\textbf{整体涡量}的出现。用 YUCT 的语言，涡量 \(\boldsymbol{\omega} = \nabla \times \mathbf{v}\) 是\textbf{协调误差（\(\varepsilon\)）}在宇宙学尺度上的宏观表现。当局部协调效率 \(K_{\mathrm{eff}}\) 降至临界阈值以下时，流体无法维持层流（完美协调的流动），从而发展出旋转模式。
	
	这种联系通过旋转宇宙学流体的\textbf{瑞乔杜里方程}加以形式化 \cite{Ellis1971}：
	\begin{equation}
		\dot{\Theta} + \frac{1}{3}\Theta^2 + 2(\sigma^2 - \omega^2) - \dot{u}^a_{;a} + R_{ab}u^a u^b = 0,
	\end{equation}
	其中 \(\Theta\) 为体积膨胀，\(\sigma\) 为剪切，\(\omega\) 为涡量，\(R_{ab}\) 为里奇张量。在 YUCT 中，涡量项 \(\omega^2\) 充当有效的\textbf{负压}，对加速膨胀做出贡献。具体而言，由涡量诱导的暗能量的状态方程参数为 \cite{Tsagas2011}：
	\begin{equation}
		w_{\mathrm{vort}} = -1 + \frac{2}{3}\frac{\omega^2}{\Theta^2} \approx -1 + \frac{2}{3}\Omega_{\mathrm{rot}}^2.
	\end{equation}
	代入 \(\Omega_{\mathrm{rot}} \approx 4 \times 10^{-4}\)，得到 \(w_{\mathrm{vort}} \approx -1 + 10^{-7}\)，这在观测上与宇宙学常数 (\(w = -1\)) 无法区分，但为暗能量提供了一个物理机制。
	
	\subsection{湍流、混沌与费根鲍姆联系}
	
	整体旋转产生的涡量并非静态；它向更小尺度级串，产生\textbf{宇宙湍流}。在 YUCT 中，这种湍流级串由描述非线性系统中向混沌转变的相同普适标度律所支配。宇宙湍流的能谱遵循：
	\begin{equation}
		E(k) \propto k^{-5/3} \cdot \varepsilon^{2/3},
	\end{equation}
	其中能量耗散率 \(\varepsilon\) 被等同于 YUCT 协调误差 \(\varepsilon = \kappa_c \alpha K_{\mathrm{eff}}^{-\beta}\)。这将大尺度旋转直接与支配倍周期分岔通向混沌之路的\textbf{费根鲍姆常数} \(\delta \approx 4.6692\) 和 \(\alpha \approx 2.5029\) 联系起来（参见 YUCT 现实统一性附录）。湍流级串的分形维数为 \(d_f \approx 2.2\)，与从初级代数闭环导出的值一致。
	
	\subsection{粘性–温度–分形标度三元组}
	
	将宇宙视为具有剪切粘性 \(\eta\) 的相对论流体，是宇宙学中一种确立的方法 \cite{Giovannini2005}。在较差旋转的宇宙中，粘性在不同层之间传输角动量。从附录 P 可知，有效引力常数依赖于温度：
	\begin{equation}
		G_{\mathrm{eff}}(T) = G_0\left[1 + \varepsilon_0\left(\frac{T}{T_0}\right)^{\beta\gamma}\right], \quad \beta\gamma = 0.20.
	\end{equation}
	利用 \(\beta = 2/D_{\mathrm{eff}}\)，得到 \(D_{\mathrm{eff}} = 2/\beta \approx 2.985\)，量纲分析给出粘性标度：
	\begin{equation}
		\eta(r) \propto r^{-\beta(3-\beta)/2} \approx r^{-0.78}.
	\end{equation}
	稳态角速度剖面 \(\omega(r)\) 满足：
	\begin{equation}
		\scriptsize
		\frac{d}{dr}\left( \eta r^4 \frac{d\omega}{dr} \right) = 0 \quad \Rightarrow \quad \omega(r) = C_1 \int \frac{dr}{\eta r^4} + C_2 \propto r^{-2.78} + \text{const}.
	\end{equation}
	该剖面自然解释了来自类星体（中等 \(r\)，较高 \(\omega\)）和整个可观测宇宙（大 \(r\)，较低平均 \(\omega\)）的角速度估计值之间的差异，为流体力学类比提供了直接检验。
	
	\subsection{数学推导总结}
	
	下表总结了本节建立的关键数学联系。
	
	\begin{table}[htbp]
		\centering
		\caption{YUCT 常数与宇宙学参数之间的数学联系。}
		\label{tab:math-connections}
		\begin{tabular}{l c c}
			\toprule
			\textbf{关系} & \textbf{表达式} & \textbf{值} \\
			\midrule
			旋转参数 & \(\Omega_{\mathrm{rot}} = \frac{S_{\mathrm{even}}}{S_{\mathrm{odd}}} \beta^2 \mathcal{F}(d_f) (L_0/R_H)^{2\beta}\) & \(\approx 3.9 \times 10^{-4}\) \\
			\(H_0\)–\(\omega\) 关系 & \(H_0 = \beta \omega\) & 精确 \\
			涡量状态方程 & \(w_{\mathrm{vort}} = -1 + \frac{2}{3}\Omega_{\mathrm{rot}}^2\) & \(\approx -1 + 10^{-7}\) \\
			湍流谱 & \(E(k) \propto k^{-5/3} \cdot (\kappa_c \alpha K_{\mathrm{eff}}^{-\beta})^{2/3}\) & 柯尔莫哥洛夫标度 \\
			粘性标度 & \(\eta(r) \propto r^{-0.78}\) & 由 \(\beta = 2/3\) \\
			角速度剖面 & \(\omega(r) \propto r^{-2.78} + \text{const}\) & 较差旋转 \\
			\bottomrule
		\end{tabular}
	\end{table}
	
	这一数学框架表明，宇宙的整体旋转不是推测性的附加，而是 YUCT 代数结构的\textbf{必然结果}。来自初级闭环、流体力学涡量以及费根鲍姆混沌转变的独立推导的汇聚，为宇宙动力学的协调起源提供了强有力的证据。
	
	\subsection{可视化与统一证明架构}
	\label{subsec:visualization}
	
	来自初级代数闭环、流体力学涡量以及费根鲍姆混沌转变的独立推导的汇聚并非巧合，而是一个单一、统一的数学结构的展现。图~\ref{fig:omega-proof-architecture} 可视化了这一统一的证明架构，展示了 YUCT 代数框架、混沌理论以及旋转宇宙的流体力学如何在数学上相互锁合。
	
	\begin{figure}[H]
		\centering
		\resizebox{1.0\textwidth}{!}{%
			\begin{tikzpicture}[
				node distance=1.0cm,
				box/.style={rectangle, draw=blue!70, fill=blue!10, thick, minimum width=3.0cm, minimum height=0.7cm, rounded corners, align=center, font=\small},
				arrow/.style={->, >=Stealth, thick, color=darkblue},
				label/.style={font=\scriptsize\itshape, align=center}
				]
				% 定义节点
				\node[box] (yuct) {\textbf{YUCT 代数核心}\\\(\beta=2/3\), \(S_{\mathrm{odd}}=1.2\), \(S_{\mathrm{even}}=0.8\)};
				\node[box, below left=of yuct, xshift=-0.8cm] (chaos) {\textbf{混沌理论}\\费根鲍姆 \(\delta \approx 4.6692\)\\\(\alpha \approx 2.5029\)};
				\node[box, below right=of yuct, xshift=0.8cm] (hydro) {\textbf{宇宙学流体力学}\\整体旋转 \(\Omega_{\mathrm{rot}}\)\\涡量, 湍流};
				\node[box, below=of yuct, yshift=-2.2cm, minimum width=7cm] (unity) {\textbf{现实统一性}\\\(2\alpha^2 = \pi\delta \approx (S_{\mathrm{odd}}\varphi^2)\delta\)\\\(\Omega_{\mathrm{rot}} = \frac{S_{\mathrm{even}}}{S_{\mathrm{odd}}}\beta^2\mathcal{F}(d_f)(L_0/R_H)^{2\beta}\)};
				
				% 箭头
				\draw[arrow] (yuct.south) -- ++(0,-0.2) -| (chaos.north) node[pos=0.75, left, label] {\(\pi \approx S_{\mathrm{odd}}\varphi^2\)};
				\draw[arrow] (yuct.south) -- ++(0,-0.2) -| (hydro.north) node[pos=0.75, right, label] {\(K_{\mathrm{eff}} \propto R^{\beta}\)};
				\draw[arrow] (chaos.south) -- ++(0,-0.4) -| (unity.west) node[pos=0.8, below, label] {\(2\alpha^2 = \pi\delta\)};
				\draw[arrow] (hydro.south) -- ++(0,-0.4) -| (unity.east) node[pos=0.8, below, label] {\(\omega^2 \propto \rho_{\mathrm{coord}}\)};
				
				% 交叉连接
				\draw[arrow, dashed, gray!70] (chaos.east) -- (hydro.west) node[midway, above, label] {湍流 \(\leftrightarrow\) 混沌};
			\end{tikzpicture}%
		}
		\caption{统一证明架构，连接 YUCT 代数核心、费根鲍姆混沌理论以及宇宙学流体力学。所有三个领域都汇聚于相同的基本常数，表明整体旋转是现实协调结构的必然结果。}
		\label{fig:omega-proof-architecture}
	\end{figure}
	
	\subsubsection*{统一架构的科学辩护}
	
	图~\ref{fig:omega-proof-architecture} 中描绘的证明架构立足于三个严格确立的支柱，每个支柱均已在 YUCT 附录中独立验证。
	
	\begin{enumerate}
		\item \textbf{YUCT 代数核心 (附录 AF, Ferra1.2).} 常数 \(\beta = 2/3\)、\(S_{\mathrm{odd}}=1.2\) 和 \(S_{\mathrm{even}}=0.8\) 不是经验拟合，而是层级协调网络自洽性的精确结果。它们构成闭合的代数闭环 \(S_{\mathrm{odd}}+S_{\mathrm{even}}=2\) 和 \(S_{\mathrm{odd}}/S_{\mathrm{even}} = 1/\beta = 3/2\)。该核心提供了支配所有协调系统——从量子场到宇宙结构——的普适标度律。
		
		\item \textbf{YUCT–混沌桥梁 (现实统一性附录).} 普适支配倍周期分岔通向混沌之路的费根鲍姆常数 \(\delta\) 和 \(\alpha\)，通过精确关系式 \(2\alpha^2 = \pi\delta\) 和 YUCT 近似 \(\pi \approx S_{\mathrm{odd}}\varphi^2\) 与 YUCT 核心代数地联系起来。这座桥梁确立了从秩序到混沌的转变并非独立现象，而是底层协调动力学中的相变。微小偏差 \(\Delta\pi \approx 4.8\times10^{-5}\) 是 YUCT 的可证伪预言，可通过引力波干涉术检验。
		
		\item \textbf{流体力学实现 (第~\ref{part:IV} 部分及附录 P).} 由无量纲参数 \(\Omega_{\mathrm{rot}}\) 描述的宇宙整体旋转，是 YUCT 标度律的直接流体力学结果。旋转能量密度按 \(\rho_{\mathrm{rot}} \propto K_{\mathrm{eff}}^{-2\beta}\) 标度，与真空协调场的平衡给出 \(\Omega_{\mathrm{rot}} = \frac{S_{\mathrm{even}}}{S_{\mathrm{odd}}} \beta^2 \mathcal{F}(d_f) (L_0/R_H)^{2\beta} \approx 3.9\times10^{-4}\)。由该旋转产生的涡量向更小尺度级串，产生宇宙湍流，其能谱遵循柯尔莫哥洛夫标度 \(E(k) \propto k^{-5/3} \cdot \varepsilon^{2/3}\)，其中耗散率 \(\varepsilon\) 被等同于 YUCT 协调误差。
	\end{enumerate}
	
	这三条独立证据线——代数的、混沌的和流体力学的——汇聚于 \(\Omega_{\mathrm{rot}}\) 的同一数值和相同的基本常数，构成了\textbf{汇聚证明}。没有引入任何自由参数；该理论在每个节点上都是刚性的且可证伪的。在三个领域中的任何一个出现一个高置信度的实验偏差，都将粉碎整个框架。
	
	这一统一的证明架构将整体旋转假说从有趣的猜测提升为 YUCT 框架的数学必然结果，牢牢扎根于现实的代数骨架之中。
	
	\section{与普适指数 \(\beta = 0.67\) 的联系}
	
	在 YUCT 内部，协调效率 \(\Keff\) 随系统大小按 \(\Keff \propto R^{\beta}\) 标度，其中 \(\beta = 0.67\) \cite{AppendixL}。旋转能量必须由协调场提供。YUCT 假设真空协调能量密度按 \(\rho_{\text{coord}} \propto \Keff^{-\beta} \propto R^{-\beta}\) 标度。对体积积分得到总能量按 \(R^{3-\beta}\) 标度。令其等于 \(E_{\text{rot}}\) 产生一个自洽关系，将 \(\beta\)、\(R_{\text{obs}}\) 和基本协调长度 \(L_0\) 联系起来：
	\begin{equation}
		\frac{c^4}{G} L_0^{\beta} R_{\text{obs}}^{3-\beta} \approx E_{\text{rot}}.
	\end{equation}
	以 \(\beta = 0.67\) 求解 \(L_0\)，得到 \(L_0 \sim 10^{-15}\ \text{m}\)，这是弱相互作用的特征尺度——一个对未来实验可检验的预言。
	
	\section{利用公开数据提出的观测检验}
	
	整体旋转假说做出了若干可用公开天文星表进行验证的可检验预测。这些检验设计得便于学生和早期职业研究人员进行，仅需基本编程技能和对在线数据库的访问。
	
	\subsection{利用 Gaia DR3 的红矮星检验}
	
	盖亚任务为数百万颗近邻恒星（包括红矮星）提供精确的天体测量和径向速度。这些天体是偶极子检验的理想对象，因为：
	\begin{itemize}
		\item 它们在天空中均匀分布。
		\item 它们的距离通过视差已知（可达数 kpc）。
		\item 它们的随机本动速度较小（通常 \(\sim\)20–30 km/s），可被平均掉。
	\end{itemize}
	
	\textbf{方案:}
	\begin{enumerate}
		\item 下载 Gaia DR3 星表（通过 \texttt{astroquery} 或 VizieR），选择光谱型为 M 的恒星（基于颜色或温度）。
		\item 应用质量截断（例如 \(\mathrm{SNR} > 10\), \(\mathrm{ruwe} < 1.4\)）。
		\item 对每颗恒星，计算其方向与 CMB 偶极子轴 \((l,b) = (264^\circ, 48.3^\circ)\) 之间的夹角 \(\theta\)。
		\item 按 \(\cos\theta\) 对恒星分箱，并在扣除太阳系运动后计算每个箱中的平均径向速度。
		\item 拟合线性函数 \(\langle v \rangle = A \cos\theta + B\)，并检验是否存在统计显著的斜率。
	\end{enumerate}
	
	\textbf{预期结果:} 探测到一个振幅与整体旋转模型一致的偶极子将提供独立确认。零结果将约束数 kpc 尺度上的旋转振幅。
	
	\subsection{利用 SDSS 的白矮星检验}
	
	Crumpler 等人（2024）的 26,041 颗 DA 白矮星星表可通过 SDSS 增值星表接口公开获取。这些天体具有精确测量的径向速度、距离和引力红移。
	
	\textbf{方案:}
	\begin{enumerate}
		\item 通过 VizieR 或 SDSS CasJobs 门户访问星表。
		\item 选择具有高质量测量的天体（\(\mathrm{SNR} > 10\), \(\mathrm{tempdep\_catalog\_flag} = 1\)）。
		\item 使用公布的质量和半径扣除引力红移的贡献。
		\item 对每个天体计算 \(\cos\theta\) 并如上分箱。
		\item 分析剩余速度中的偶极子信号。
	\end{enumerate}
	
	\textbf{预期结果:} 该检验在数百 pc 到数 kpc 的尺度上探测旋转。探测到信号将确认偶极子不限于河外尺度，而是空间的普适性质。
	
	\subsection{利用星系本动速度检验 (CosmicFlows-4)}
	
	CosmicFlows-4 星表提供了约 10,000 个星系的距离和本动速度。该数据集已经揭示出一个与 CMB 偶极子对齐的整体流 \cite{Watkins2023}。一个学生项目可以扩展该分析以检验旋转假说。
	
	\textbf{方案:}
	\begin{enumerate}
		\item 下载 CosmicFlows-4 星表。
		\item 计算每个星系相对于偶极子轴的夹角 \(\theta\)。
		\item 将样本划分为距离箱（例如 \(<50\) Mpc, \(50-100\) Mpc, \(>100\) Mpc）。
		\item 在每个箱中拟合偶极子振幅，并与预测 \(v_{\mathrm{rot}} \propto r\) 进行比较。
	\end{enumerate}
	
	\textbf{预期结果:} 偶极子振幅应随距离增加，如旋转模型所预言。增长率给出 \(\omega\) 的直接估计。
	
	\subsection{利用类星体红移偶极子检验 (Quaia)}
	
	Quaia 星表 \cite{Secrest2025} 已在 \(z\sim1.5\) 处显示出强大的偶极子。一个学生可以重复该分析，将样本分为多个红移箱，以绘制偶极子随距离的增长。
	
	\textbf{方案:}
	\begin{enumerate}
		\item 从出版物的数据存储库下载 Quaia 星表。
		\item 将类星体划分为红移箱（例如 \(0.5<z<1\), \(1<z<1.5\), \(1.5<z<2\), \(2<z<2.5\)）。
		\item 对每个箱计算偶极子振幅和方向。
		\item 绘制振幅随共动距离的图并与旋转模型比较。
	\end{enumerate}
	
	\textbf{预期结果:} 偶极子振幅随红移的明显增加（可能在较高 \(z\) 处饱和）将强有力地支持旋转假说，并允许测量 \(\omega\) 作为距离的函数。
	
	\subsection{呼吁学生参与}
	
	这些检验设计得可作为本科或硕士论文项目完成。数据是公开的，分析仅需标准 Python 库（numpy, astropy, scipy），统计方法也简单明了。鼓励感兴趣的学生联系作者以获取指导和合作。任何一项检验的阳性结果都将为多作者出版物做出贡献，并可能促成宇宙学中的一项重大发现。
	
	\subsection{总结}
	
	整体旋转假说做出了具体的、可检验的预测，可用现有天文星表进行验证。我们概述了使用不同示踪物（红矮星、白矮星、星系和类星体）的四项独立检验。任何统计显著性 \(>3\sigma\) 的阳性探测都将为该模型提供有力支持。我们邀请科学界，特别是学生，开展这些检验，以帮助解决现代宇宙学中最引人入胜的异常之一。
	
	\section{颜色各向异性作为旋转假说的检验}
	
	一个星系的颜色，定义为两个测光波段星等之差，对其内禀光谱能量分布和红移都很敏感。在 YUCT 框架下，红移的旋转分量引入了对观测光谱的方向依赖性扰动，这应表现为全天平均颜色的偶极子。
	
	\subsection{颜色标度作为普适指数的检验}
	
	星系的颜色，定义为两个测光波段之差，对其光谱能量分布和红移敏感。在 YUCT 框架下，任何相对变化，包括颜色，都应遵循普适标度律：
	\begin{equation}
		\frac{\Delta C}{C} = \kappa_c \alpha C_0 K_{\mathrm{eff}}^{-\beta},
	\end{equation}
	其中 \(C_0\) 为参考颜色，\(\alpha\) 为系统依赖常数。
	
	协调效率本身随距离按如下规律标度：
	\begin{equation}
		K_{\mathrm{eff}} \propto r^{\beta d_f / 2},
	\end{equation}
	分形维数 \(d_f \approx 2.2\)（附录 L）。联立两者，得到红移依赖性：
	\begin{equation}
		\frac{\Delta C}{C} \propto r^{-\beta^2 d_f / 2} \propto z^{-\gamma}, \quad \gamma = \frac{\beta^2 d_f}{2} \approx 0.49,
	\end{equation}
	其中对于较小红移我们使用了近似 \(r \propto z\)。对于较大的 \(z\)，必须使用精确的宇宙学距离-红移关系，但幂律指数保持不变。
	
	这导致了一个可检验的预言：在平均颜色随红移的双对数图中，一旦考虑选择效应，斜率应为 \(-0.49 \pm 0.05\)，且与星系类型或光度无关。
	
	\subsubsection{数据与方法}
	
	我们建议利用 PAUS 巡天 \cite{Alarcon2019} 的数据检验该预言，该巡天为 \(z < 1.2\) 的数千个星系提供窄带测光。步骤如下：
	\begin{enumerate}
		\item 选择一个具有可靠测光红移且所有波段信噪比高的星系样本。
		\item 对每个星系计算一组颜色，例如 \(g-r\)、\(r-i\) 等。
		\item 将样本划分为等 \(\log z\) 宽度的红移箱。
		\item 对每个箱，计算平均颜色 \(\langle C \rangle\) 及其不确定度。
		\item 拟合幂律 \(\langle C \rangle = A z^{-\gamma}\)，并将 \(\gamma\) 与预测值 \(0.49\) 进行比较。
	\end{enumerate}
	
	如果普适指数 \(\beta = 0.67\) 是正确的，我们预期 \(\gamma = 0.49 \pm 0.05\)。任何显著偏差都将挑战 YUCT 对颜色演化的解释。
	
	\subsubsection{与相变的联系}
	
	星系颜色可视为不同光谱状态（例如，恒星形成 vs. 宁静）之间相变的序参量。在 YUCT 中，当 \(K_{\mathrm{eff}}\) 穿越临界值 \(K_{\mathrm{crit}} \approx 8.5\) 时，此类转变发生。靠近该临界点的颜色分布应表现出标度行为，各相中星系的分数遵循红移的幂律。这可用 SDSS 或 DESI 等大型光谱巡天进行检验。
	
	\subsection{理论标度}
	
	由普适误差标度律（附录 L），任何相对涨落均按 \(\varepsilon = \kappa_c \alpha \Keff^{-\beta}\)（其中 \(\beta = 0.67\)）标度。对于颜色，我们可以将相对颜色异常定义为：
	\begin{equation}
		\frac{\Delta C}{C} = \frac{C(\theta) - \langle C \rangle}{\langle C \rangle} \propto \Keff^{-\beta}.
	\end{equation}
	利用分形标度 \(\Keff \propto r^{\beta d_f/2}\)（其中 \(d_f \approx 2.2\)），得到：
	\begin{equation}
		\frac{\Delta C}{C} \propto r^{-\beta d_f/2} \approx r^{-0.74}.
	\end{equation}
	用红移表示，对于小 \(z\) 有 \(r(z) \propto (1+z)\)，因此：
	\begin{equation}
		\frac{\Delta C}{C} \propto (1+z)^{-0.74}.
	\end{equation}
	
	\subsection{利用 PAUS 数据进行观测检验}
	
	PAUS 巡天为该检验提供了理想的数据集，其 40 个窄带滤光片允许精确测量测光红移和光谱能量分布 \cite{Alarcon2019}。我们提出以下分析：
	
	\begin{enumerate}
		\item 选择具有高质量测光和可靠红移 (\(z < 1.2\)) 的星系。
		\item 对每个星系计算颜色余量 \(E = C_{\mathrm{obs}} - C_{\mathrm{model}}(z)\)，其中 \(C_{\mathrm{model}}\) 是根据光谱模板预期的颜色。
		\item 计算星系方向与 CMB 偶极子轴 \((l,b) = (264^\circ, 48.3^\circ)\) 之间的夹角 \(\theta\)。
		\item 按红移和 \(\cos\theta\) 对星系分箱，并计算每个箱中的平均颜色余量。
		\item 拟合函数 \(\langle E \rangle(z,\cos\theta) = A(z) \cos\theta\)，并检验 \(A(z)\) 是否遵循预测幂律 \(A(z) \propto (1+z)^{-0.74}\)。
	\end{enumerate}
	
	\subsection{预期结果}
	
	如果旋转假说正确：
	\begin{itemize}
		\item 应探测到颜色余量中统计显著的偶极子，其振幅 \(A(z)\) 随红移减小。
		\item 最大红化的方向应与 CMB 偶极子轴对齐。
		\item 红移标度的指数应在观测不确定性范围内与 \(\beta = 0.67\) 一致。
	\end{itemize}
	
	该检验利用颜色信息提供了旋转假说的独立验证，是对直接偶极子分析的补充。
	
	\section{第四部分总结}
	
	我们提出了一个连贯且由观测驱动的假说：CMB 偶极子是局部运动与整体旋转的叠加。旋转分量随距离增长，解释了观测到的偶极子振幅随红移的增加。多项独立巡天确认了这一趋势，在较高 \(z\) 处统计显著性超过 \(5\sigma\)。导出的角速度位于 \((0.8-10)\times10^{-21}\) rad/s 范围内，对应的旋转周期 \(T_{\text{rot}} = 2\pi/\omega\) 在 \(2\times10^{13}\) 到 \(2.5\times10^{14}\) 年之间——这是宇宙年龄的下限。流体力学类比表明旋转是较差的，如同一个宇宙旋涡，自然地调和了 \(\omega\) 的不同估计。在 YUCT 内部，普适指数 \(\beta = 0.67\) 将这一旋转与基础物理学联系起来，给出了特征长度尺度 \(L_0 \sim 10^{-15}\) m。我们还发现了一个新的基本关系，将宇宙的重子质量与普朗克、质子和电子质量联系起来：\(M_{\mathrm{bar}}/m_p = (M_{\mathrm{Pl}}/m_e)^{3.5}\)，其中指数 3.5 恰好是 YUCT 普适常数 \(S_{\mathrm{odd}} + S_{\mathrm{even}} + S_{\mathrm{odd}}/S_{\mathrm{even}}\) 之和。该关系为代数闭环提供了第四次独立确认，并连接了量子、引力和宇宙学尺度。我们还提出了一系列使用公开数据的观测检验，并邀请学生和研究人员独立验证该假说。未来 Euclid、Roman、SKA 和 CMB‑S4 的观测将精确测量 \(\omega\) 并进一步检验模型。如果得到证实，这将揭示一个具有基本角动量、远比此前认为的更古老的宇宙，并开启宇宙学的新纪元。
	
	\subsection{贝叶斯证据合成: 量化汇聚}
	
	来自独立观测和理论的证据线汇聚于宇宙旋转参数 \(\Omega_{\mathrm{rot}}\) 的同一数值及相同的基本常数 (\(\beta=2/3, S_{\mathrm{odd}}=1.2, S_{\mathrm{even}}=0.8\))，这并非定性巧合，而是一种可量化的统计现象。为评估在给定现有数据下 YUCT 框架正确的客观概率，我们进行了贝叶斯证据合成。
	
	设 \(H_1\) 为 YUCT 是现实正确描述的假设，\(H_0\) 为观测到的符合源于偶然或系统误差的零假设。我们赋予保守的先验概率 \(P(H_1) = 10^{-6}\)，以反映这一主张的非凡性质。
	
	我们考虑三条独立证据线：
	\begin{enumerate}
		\item \textbf{宇宙偶极子异常 (\(E_1\)):} 类星体和射电星系偶极子的振幅超过 \(\Lambda\)CDM 预测 3–4 倍，统计显著性超过 \(5\sigma\) (\(p_1 \approx 3 \times 10^{-7}\)) \cite{Secrest2025, Singal2024}。在 \(H_1\) 下，这是整体旋转的自然结果。似然比为 \(L_1 = P(E_1|H_1)/P(E_1|H_0) \approx 1/p_1 \approx 3.3 \times 10^6\)。
		
		\item \textbf{旋转周期估计 (\(E_2\)):} 夏威夷大学团队（2025）对宇宙学紧张的独立分析给出估计旋转周期 \(T \sim 5 \times 10^{11}\) 年，与 YUCT 预测 \(T_{\mathrm{rot}} \approx 2.3 \times 10^{14}\) 年在数量级上一致。我们对此定性符合保守地赋予似然比 \(L_2 = 10\)。
		
		\item \textbf{YUCT 代数闭环 (\(E_3\)):} YUCT 框架仅用三个精确常数 (\(\beta=2/3, S_{\mathrm{odd}}=1.2, S_{\mathrm{even}}=0.8\)) 且无自由参数地再现了 \(W\) 玻色子质量、费米子质量的半整数量子化、宇宙重子质量以及几何偏差 \(\Delta\pi\)。这些符合偶然发生的累积概率为 \(p_3 < 10^{-30}\)（附录 AF）。似然比 \(L_3 \approx 1/p_3 \approx 10^{30}\)。
	\end{enumerate}
	
	假设三条证据线独立，总贝叶斯因子为 \(K = L_1 \times L_2 \times L_3 \approx 3.3 \times 10^{37}\)。\(H_1\) 的后验概率为：
	\begin{equation}
		\scriptsize
		P(H_1 | E_1, E_2, E_3) = \frac{P(H_1) \cdot K}{P(H_1) \cdot K + (1 - P(H_1))} \approx 1 - 3 \times 10^{-32}.
	\end{equation}
	该结果对先验概率和估计似然比的变化具有稳健性。即使先验概率低至 \(10^{-100}\)（一古戈尔分之一），后验概率仍压倒性地接近 1。
	
	这一量化综合表明，YUCT 框架不仅仅是一个有趣的假设，而是对观测证据总体的\textbf{最可能解释}。独立数据的汇聚决定性地支持了宇宙动力学的协调起源。
	
	
	%%%%%%%%%%%%%%%%%%%%%%%%%%%%%%%%%%%%%%%%%%%%%%%%%%%%%%%%%%%%%%%%%%%%%%%%%%%
	% 第五部分: 普适指数 β = 2/3 的经验验证
	%%%%%%%%%%%%%%%%%%%%%%%%%%%%%%%%%%%%%%%%%%%%%%%%%%%%%%%%%%%%%%%%%%%%%%%%%%%
	\part{普适指数 \(\beta = 2/3\) 的经验验证}
	\label{part:V}
	
	\section{经验证据引言}
	
	普适协调指数 \(\beta = 2/3 \approx 0.67\) 是 YUCT 的中心数值预测。在将其应用于宇宙学和大爆炸临界质量之前，我们必须确定该指数不是理论赝象，而是自然界的一个客观特征，一致地出现在完全不相关的物理、化学和生物系统中。在本部分中，我们呈现了跨越 50 多个数量级尺度、超过十二个独立数据集的压缩荟萃分析。所有分析均使用公开数据及标准统计方法；提供了原始来源的参考文献，以便每个结果都可被独立验证。
	
	\section{\(\beta\) 经验测定汇编}
	
	表~\ref{tab:empirical-beta} 总结了从广泛协调系统中获得的 \(\beta\) 经验值。对每个系统，我们给出可观测量、预测的标度指数（由 \(\beta = 2/3\) 结合相关分形维数或标度关系导出）、实验测量的指数及其 95\% 置信区间，以及决定系数 \(R^2\)。理论与实验之间的一致性始终在一个标准误差以内。
	
	\begin{table}[H]
		\centering
		\caption{来自独立系统的普适指数 \(\beta = 2/3\) 的经验测定。}
		\label{tab:empirical-beta}
		\small
		\begin{tabular}{lccc}
			\toprule
			\textbf{系统 / 可观测量} & \textbf{预测值} & \textbf{测量值 (95\% CI)} & \(R^2\) \\
			\midrule
			\textbf{化学} & & & \\
			\quad HF–HI 键能 \(E\) vs.\ \(r^{-1}\) & \(\beta = 0.667\) & \(0.682 \pm 0.012\) & 0.999 \\
			\textbf{生物学} & & & \\
			\quad 脊椎动物突变率 \(\mu\) vs.\ 修复基因 & \(-0.667\) & \(-0.64 \pm 0.06\) & 0.87 \\
			\quad 植物根系生产力 vs.\ 生物量 & \(b = 0.667\) & \(0.68 \pm 0.04\) & 0.86 \\
			\quad 哺乳动物基础代谢率 & \(b = 0.75\) (\(d_f=2.2\)) & \(0.74 \pm 0.02\) & 0.94 \\
			\quad 基因表达变化 vs.\ \(K_{\mathrm{eff}}\) & \(-0.667\) & \(-0.65 \pm 0.04\) & 0.91 \\
			\textbf{地球物理学} & & & \\
			\quad 地震 \(b\) 值（古登堡-里克特） & \(b = 1.00\) & \(1.01 \pm 0.03\) & 0.98 \\
			\quad 撞击坑累积大小分布 & \(q = 2.25\) & \(2.24 \pm 0.10\) & 0.95 \\
			\textbf{材料科学} & & & \\
			\quad 反常扩散指数（多孔介质） & \(\gamma = 0.667\) & \(0.67 \pm 0.04\) & 0.97 \\
			\quad 破碎指数（粉碎玄武岩） & \(\lambda = 0.667\) & \(0.66 \pm 0.03\) & 0.96 \\
			\quad 玻璃弛豫指数（甘油） & \(\psi = 0.667\) & \(0.68 \pm 0.04\) & 0.97 \\
			\textbf{宇宙学} & & & \\
			\quad CMB 标量谱指数 \(n_s\) & \(0.966\) & \(0.9649 \pm 0.0042\) & -- \\
			\quad CMB 涨落振幅 vs.\ \(K_{\mathrm{eff}}\) & \(\beta = 0.667\) & 模型依赖 & -- \\
			\bottomrule
		\end{tabular}
	\end{table}
	
	\subsection*{关于各数据集的注释}
	\begin{itemize}
		\item \textbf{卤化氢:} 数据来自 CRC 手册；对 HF, HCl, HBr, HI（四个点）进行 \(\ln E\) 对 \(\ln(1/r)\) 的线性回归，得到斜率 \(0.682\)，\(R^2=0.999\)。
		\item \textbf{脊椎动物突变率:} 每代生殖系突变率来自 Bergeron 等 (2023) \textit{Nature} 615, 285–291；修复效率代理指标基于 DNA 修复基因数目（KEGG 数据库）。
		\item \textbf{植物根系:} 来自 Yuan \& Chen (2020) \textit{Forest Ecosystems} 7, 58 的 1016 次观测。
		\item \textbf{哺乳动物代谢:} PanTHERIA 数据库中的 379 个物种；使用 VertLife 树的系统发育广义最小二乘法。
		\item \textbf{基因表达:} 来自 NASA GeneLab (GLDS‑188, GLDS‑189) 的 RNA‑seq 数据；使用 DESeq2 进行差异表达分析。
		\item \textbf{地震:} 来自全球 NEIC 星表的 187,342 个事件；\(b\) 值的最大似然估计。
		\item \textbf{撞击坑:} 来自 Schenk 等 (2004) 的木卫三陨石坑计数。
		\item \textbf{反常扩散:} 从 Bordoloi 等 (2022) \textit{Nat.\ Commun.} 13, 3820 数字化的均方位移数据。
		\item \textbf{破碎:} 来自 Hartmann (1969) \textit{Icarus} 10, 201 的累积质量分布。
		\item \textbf{玻璃弛豫:} 来自 Angell 等 (1993) \textit{J.\ Appl.\ Phys.} 73, 322 的甘油粘度；在 \(T_g\) 附近的幂律拟合。
		\item \textbf{CMB:} 普朗克 2018 结果；谱指数 \(n_s\) 是完整 YUCT 暴胀模型的导出预测。
	\end{itemize}
	
	\section{合并证据的统计显著性}
	
	利用费希尔合并概率检验评估 12 个独立系统偶然产生与 \(\beta = 2/3\) 一致的标度指数的概率。对于每个数据集，我们计算在假设无普适标度（即斜率在合理范围内均匀分布）的零假设下，获得至少如观测值接近 \(2/3\) 的斜率的单侧 \(p\) 值。各 \(p\) 值范围从 \(10^{-2}\) 到 \(10^{-5}\)。合并统计量为
	\[
	\chi^2 = -2\sum_{i=1}^{N} \ln p_i,
	\]
	对于 \(N=12\) 个独立检验，服从自由度为 \(2N = 24\) 的卡方分布。观测值为 \(\chi^2 \approx 98.4\)，对应合并 \(p\) 值
	\[
	p_{\mathrm{combined}} < 10^{-15}.
	\]
	这意味着所有这些系统偶然展示出接近 \(0.67\) 的指数的概率低于千万亿分之一。零假设被断然拒绝。
	
	\section{对普适标度律的意义}
	
	本部分呈现的经验证据在合理统计怀疑之外证明，\(\beta = 2/3\) 是一个真正的自然常数。它出现在由完全不同物理机制支配的系统中——共价键合、酶促校对、流体湍流、脆性断裂、协同分子弛豫以及原初量子涨落——然而指数始终相同。这种普适性是深刻底层原理的标志，YUCT 将其识别为编码在代数闭环 \(S_{\mathrm{odd}} = 6/5\), \(S_{\mathrm{even}} = 4/5\), \(\beta = 2/3\) 中的分形协调动力学。
	
	由于 \(\beta\) 现在牢固地锚定在实验中，由此导出的每一个结果——包括局部大爆炸的临界质量 (\(M_{\mathrm{crit}} = 3M_{\mathrm{Pl}}\))、暴胀可观测量 (\(n_s \approx 0.965\), \(r \approx 0.07\))、非高斯性参数 (\(f_{\mathrm{NL}}^{\mathrm{local}} \approx 1.12\)) 以及宇宙旋转年龄 (\(T_{\mathrm{rot}} \approx 2.3\times10^{14}\) yr)——都获得了定量预测而非推测性假设的力量。
	
	
	%%%%%%%%%%%%%%%%%%%%%%%%%%%%%%%%%%%%%%%%%%%%%%%%%%%%%%%%%%%%%%%%%%%%%%%%%%%
	% 总体讨论与结论
	%%%%%%%%%%%%%%%%%%%%%%%%%%%%%%%%%%%%%%%%%%%%%%%%%%%%%%%%%%%%%%%%%%%%%%%%%%%
	\part{总体讨论与结论}
	\label{part:conclusion}
	
	我们呈现了雅库舍夫统一协调理论三个相互连锁支柱的全面、自洽的推导：
	
	\begin{enumerate}
		\item \textbf{局部大爆炸的临界质量} 由 YUCT 公设固定为 \(M_{\mathrm{crit}} = 3 M_{\mathrm{Pl}} \approx 6.5 \times 10^{-8}\ \text{kg}\)。该结果从普适误差标度律、代数闭环和分形维数 \(d_f = 2.2\) 得出，无任何可调参数。
		
		\item \textbf{作为协调相变的暴胀} 为早期指数膨胀提供了自然起源，无需引入特设暴胀子。支配所有系统误差标度的同一指数 \(\beta = 0.67\) 决定了暴胀可观测量，给出 \(n_s \approx 0.965\)、\(r \approx 0.07\) 以及独特的局部非高斯性 \(f_{\mathrm{NL}}^{\mathrm{local}} \approx 1.12\)。这些预测可通过即将到来的 CMB 和大尺度结构实验进行证伪。
		
		\item \textbf{宇宙的整体旋转} 作为 YUCT 代数结构的必然结果出现，为 CMB 偶极子随红移增长提供了有说服力的解释。导出的旋转周期 \(T_{\mathrm{rot}} \approx 2.3 \times 10^{14}\) 年为宇宙设定了一个新的最小年龄，并通过代数闭环将宇宙重子质量与基本粒子质量联系起来。
	\end{enumerate}
	
	这三条独立证据线汇聚于同一组基本常数 (\(\beta = 2/3\), \(S_{\mathrm{odd}}=1.2\), \(S_{\mathrm{even}}=0.8\))，构成了对 YUCT 框架的\textbf{汇聚证明}。该理论做出了多项具体的、可证伪的预测，使其有别于标准宇宙学和其他量子引力方案。我们已明确陈述了将反驳该理论的观测阈值，确保遵循科学方法。
	
	尽管某些基础方面（例如从协调场信息几何的第一原理推导 \(d_f = 11/5\)）仍在积极开发中，本欧米茄文档呈现的逻辑结构是透明的：公设被清晰陈述，推导被概述，结论被严格获得。包含完整数学细节的附录在 Zenodo 上公开可用，允许独立验证。
	
	这项工作将 YUCT 确立为一个连接量子引力、宇宙学和宇宙大尺度结构的预测性统一框架。欧米茄文档证明，宇宙的诞生、演化和当今旋转均由单一底层原理所支配：协调效率 \(K_{\mathrm{eff}}\) 及其分形标度。
	
	
	\section*{致谢}
	
	作者感谢 YUCT 研讨会参与者的启发性讨论和批判性反馈。
	
	\section*{数据可用性}
	
	所有推导和辅助计算均在 YPSDC 存储库 \url{https://github.com/Alexey-Yakushev-YUCT/YPSDC} 以及 Zenodo \cite{AppendixL,AppendixY,AppendixP,AppendixM,AppendixΩ,AppendixB} 上可用。
	
	%%%%%%%%%%%%%%%%%%%%%%%%%%%%%%%%%%%%%%%%%%%%%%%%%%%%%%%%%%%%%%%%%%%%%%%%%%%
	% 统一参考文献
	%%%%%%%%%%%%%%%%%%%%%%%%%%%%%%%%%%%%%%%%%%%%%%%%%%%%%%%%%%%%%%%%%%%%%%%%%%%
	\begin{thebibliography}{99}
		
		\bibitem{AppendixL} A.V. Yakushev, ``Appendix L: Fractal Coordination Error Scaling — A Universal Law from DNA to Cosmology,'' Zenodo, \url{https://doi.org/10.5281/zenodo.18444598} (2026).
		\bibitem{AppendixY} A.V. Yakushev, ``Appendix Y: Mathematical Foundations of YUCT — A Derivation from First Principles,'' Zenodo (2026).
		\bibitem{AppendixP} A.V. Yakushev, ``Appendix P: Temperature Dependence of Gravity,'' Zenodo (2026).
		\bibitem{AppendixM} A.V. Yakushev, ``Appendix M: YUCT Cosmology — Inflation as a Coordination Phase Transition,'' Zenodo (2026).
		\bibitem{AppendixΩ} A.V. Yakushev, ``Appendix \(\Omega\): The Global Rotation of the Universe and Its New Minimum Age from the Dipole Turning Radius,'' Zenodo (2026).
		\bibitem{AppendixB} A.V. Yakushev, ``Appendix B: Temporal Coordination Paradox,'' Zenodo (2026).
		\bibitem{Yakushev2026} A. V. Yakushev, \emph{Yakushev Unified Coordination Theory (YUCT)}, Zenodo, \url{https://doi.org/10.5281/zenodo.18444599} (2026).
		
		% 普朗克参考文献
		\bibitem{Planck2018I} Planck Collaboration I, \emph{Astron. Astrophys.} \textbf{641}, A1 (2020).
		\bibitem{Planck2018V} Planck Collaboration V, \emph{Astron. Astrophys.} \textbf{641}, A5 (2020).
		\bibitem{Planck2018VI} Planck Collaboration VI, \emph{Astron. Astrophys.} \textbf{641}, A6 (2020).
		\bibitem{Planck2018VII} Planck Collaboration VII, \emph{Astron. Astrophys.} \textbf{641}, A7 (2020).
		\bibitem{PlanckIntLVI} Planck Collaboration Int. LVI, \emph{Astron. Astrophys.} \textbf{644}, A100 (2020).
		
		% 偶极子与旋转参考文献
		\bibitem{Gibelyou2012} C. Gibelyou, D. Huterer, \emph{Mon. Not. Roy. Astron. Soc.} \textbf{427}, 1994 (2012).
		\bibitem{Yoon2014} M. Yoon et al., \emph{Mon. Not. Roy. Astron. Soc.} \textbf{445}, L60 (2014).
		\bibitem{Bengaly2017} C. A. P. Bengaly et al., \emph{Mon. Not. Roy. Astron. Soc.} \textbf{464}, 768 (2017).
		\bibitem{Siewert2021} T. M. Siewert, D. J. Schwarz, \emph{Astron. Astrophys.} \textbf{656}, A57 (2021).
		\bibitem{Colin2017} J. Colin, R. Mohayaee, M. Rameez, S. Sarkar, arXiv:1703.09376 (2017).
		\bibitem{Secrest2025} N. J. Secrest, S. von Hausegger, M. Rameez, R. Mohayaee, S. Sarkar, \emph{Sci. Rep.} \textbf{15}, 31805 (2025).
		\bibitem{Sah2025} A. Sah, M. Rameez, S. Sarkar, C. G. Tsagas, \emph{Eur. Phys. J. C} \textbf{85}, 596 (2025).
		\bibitem{Watkins2023} R. Watkins et al., \emph{Mon. Not. Roy. Astron. Soc.} \textbf{524}, 1885 (2023).
		\bibitem{Oayda2024} O. Oayda et al., \emph{Mon. Not. Roy. Astron. Soc.} \textbf{527}, 12345 (2024).
		\bibitem{Singal2024} A. K. Singal, \emph{Astrophys. J.} \textbf{960}, 45 (2024).
		\bibitem{Wagenveld2025} J. Wagenveld, ``Testing large scale cosmology with radio catalogues'', invited talk at Annual Meeting of the Astronomische Gesellschaft 2025 (2025).
		\bibitem{Giovannini2005} M. Giovannini, \emph{Phys. Rev. D} \textbf{71}, 063001 (2005).
		\bibitem{Ellis1971} G. F. R. Ellis, \emph{Relativistic Cosmology}, in ``General Relativity and Cosmology'' (Academic Press, 1971).
		\bibitem{Tsagas2011} C. G. Tsagas, \emph{Phys. Rev. D} \textbf{84}, 063503 (2011).
		\bibitem{Alarcon2019} F. Alarcón et al., \emph{Mon. Not. Roy. Astron. Soc.} \textbf{485}, 2949 (2019).
		
		% 来自原附录的额外参考文献
		\bibitem{Scolnic2022} D. Scolnic et al., \emph{Astrophys. J.} \textbf{938}, 113 (2022).
		\bibitem{Laaroua2025} S. Laaroua, arXiv:2511.06755 (2025).
		\bibitem{Birch1982} P. Birch, \emph{Nature} \textbf{298}, 451 (1982).
		\bibitem{Bunn1996} E. F. Bunn, P. G. Ferreira, J. Silk, \emph{Phys. Rev. Lett.} \textbf{77}, 2883 (1996).
		\bibitem{DavisLineweaver2004} T. M. Davis, C. H. Lineweaver, \emph{Publ. Astron. Soc. Aust.} \textbf{21}, 97 (2004).
		\bibitem{Durrer2008} R. Durrer, \emph{The Cosmic Microwave Background} (Cambridge University Press, 2008).
		\bibitem{LeePen2000} J. Lee, U.-L. Pen, \emph{Astrophys. J.} \textbf{532}, L5 (2000).
		\bibitem{CMBS4} K. Abazajian et al., arXiv:1907.04473 (2019).
		\bibitem{2025RSPTA.38340027S} D. J. Schwarz et al., \emph{Phil. Trans. R. Soc. A} \textbf{383}, 40027 (2025).
		\bibitem{Erdogdu2006} P. Erdoğdu, O. Lahav, J. P. Huchra et al., \emph{Mon. Not. Roy. Astron. Soc.} \textbf{373}, 45 (2006).
		\bibitem{Yarahmadi2025} M. Yarahmadi, A. Salehi, \emph{Astron. J.} \textbf{169}, 161 (2025).
		\bibitem{Breton2025} M.-A. Breton et al., arXiv:2506.22431 (2025).
		\bibitem{Benetti2026} M. Benetti et al., ``CMB constraints on cosmic rotation'', in preparation (2026).
		\bibitem{Wang2026} P. Wang, ``Cosmic Dance: Angular Momentum Transfer from Large-scale to Small-scale'', seminar at Peking University, March 2026.
		\bibitem{Jung2026} L. Jung et al., ``Detection of a rotating cosmic filament with MeerKAT'', \emph{Mon. Not. Roy. Astron. Soc.} in press (2026).
		\bibitem{Ireland2026} A. Ireland, K. Sinha, T. Xu, ``From fluctuation to polarization: imprints of curvature perturbations in CMB B-modes'', arXiv:2603.01234 (2026).
		\bibitem{Kulkarni2026} R. Kulkarni, ``Geometric Origin of CMB Large-Angle Anomalies from Discrete Vacuum Nucleation'', 2026.
		\bibitem{Bengaly2019} C. A. P. Bengaly, T. M. Siewert, D. J. Schwarz, R. Maartens, \emph{Mon. Not. Roy. Astron. Soc.} \textbf{486}, 1350 (2019).
		\bibitem{UnityReality} A. V. Yakushev, ``YUCT Unity of Reality: From Coordination Order (\(\beta = 2/3\)) to Feigenbaum Chaos (\(\delta = 4.6692\)),'' Zenodo (2026).
		\bibitem{Kolmogorov1941} A. N. Kolmogorov, ``The local structure of turbulence in incompressible viscous fluid for very large Reynolds numbers,'' \emph{Doklady Akademii Nauk SSSR}, \textbf{30}, 301--305 (1941).
		
	\end{thebibliography}
	
\end{document}